\documentclass[11pt,fleqn]{report}
\usepackage{graphicx}
\usepackage{amssymb}
\usepackage{longtable}
\usepackage[T1]{fontenc}

\usepackage{cml}
\lstnewenvironment{codeXML}
    {\lstset{basicstyle=\small\ttfamily,
             language=XML,
             xleftmargin=0.0in,
             xrightmargin=0.0in,
             backgroundcolor=\color{Orchid!60}, % 60% orchid, 40% white
             escapeinside={\%*}{*)}, %  .tex file characters inside %*...*) are
                                     %  treated as Tex commands.
            }
    }
    { }

\setcounter{secnumdepth}{4}

\newcommand{\ModelName}{Fault Management}
\newcommand{\ModelAuthor}{Daniel Ghan \\ Gary Turner}
\newcommand{\ModelAffiliation}{Odyssey Space Research}
\newcommand{\ModelDate}{April 2022}

\pagestyle{plain}
\usepackage[parfill]{parskip}
\usepackage[margin=1.5cm]{caption}
\usepackage{subfig}
\usepackage{enumitem}
\setlist[itemize]{noitemsep}


\begin{document}
\titlePage
\clearpage % Reset page number
\pagenumbering{roman}
\tableofcontents % Print table of contents
\vfill
{\Large \textbf {Revision History}}
{\large
\begin{center}
\begin{tabular}{|l||l|l|l|}
\hline
 \textbf{Version} & \textbf{Date} & \textbf{Author} & \textbf{Purpose} \\ \hline
 \hline
1 & April 2022 & Daniel Ghan & Initial Version \\ \hline
2 & September 2022 & Gary Turner & Extended Description and Verification \\ \hline
\end{tabular}
\end{center}
}

\chapter{Introduction} % Make a "chapter" heading
\pagenumbering{arabic} % Start text with arabic 1
\par The Fault Management model is designed primarily to support injecting
faults into simulated sensor readings; it may also be used for faulting
other models, including effectors.
The faults and their trigger conditions can be instantiated either at compile
time (directly in simulation objects) or at run time (by instantiating a fault
manager and defining the faults and triggers in an XML file).

\par This chapter summarizes the features and limitations of the Fault
Management model; details on how to implement it can be found in the User's
Guide (Chapter \ref{sec:UG}).

\section{Conceptual Elements} \label{intro_concept}
\subsection {Fault}
\par A Fault (see Section \ref{sec:intro:faults}) is some modification that is
applied to the value of a specified
simulation variable. In general, the Fault Management model supports ongoing
modification of the specified variable with a value that may evolve in time.
While a Fault may be used to apply an isolated, single modification (e.g. flip a
boolean flag), such a capability may also be implemented as an event (see CML
events-manager model), or a discrete event (see CML trajectory model). The
strength of the Fault Management capability lies in being able to consistently
apply values to, for example, apply white noise, or a walk-off error, etc. to
the evolving state of some variable.
\subsection{Trigger}
\par A trigger is a logical condition that \textbf{can} lead to the
application of a fault. The connection between triggers and Faults is
not direct; triggers and Faults are connected through the concept of a
trigger-group, explored in the next section.

A trigger typically comprises a \textit{trigger-variable}, a
\textit{trigger-value}, and a logical comparator.
The trigger is considered \textit{triggered} when the value of the
\textit{trigger-variable} satisfies the specified logical comparison with the
\textit{trigger-value}.
For example, a trigger set to compare the vehicle altitude against some
specified height (e.g. some ceiling) may be used to initiate a Fault that
applies noise to a signal when the
vehicle passes below that threshold height.

A trigger also has an \textit{enabled} property. A trigger that is not enabled
plays no role in determining whether a fault is applied. An enabled trigger
\textbf{may} influence the application of a Fault via a trigger-group.

\subsection{Trigger Group}
\par As may be expected from the name, a trigger-group is a collection of
(one or more) triggers. Triggers are collected into a trigger-group using the
boolean AND logic:
a \textbf{trigger group is \textit{triggered}} when ALL of its \textit{enabled}
triggers are \textit{triggered}. Note -- the \textit{triggered} status of any
non-enabled triggers in the group are not considered.

Each Fault uses one or more trigger-groups, collected using the boolean OR
logic: a \textbf{Fault is \textit{triggered}} when ANY of its trigger-groups are
\textit{triggered}.

A \textit{triggered} Fault will apply its modification to
the target variable.


\section{Faults} \label{sec:intro:faults}
\par Each Fault consists of a reference to a simulation variable, at least one
trigger group, and some parameters specific to the nature of the fault being
applied.
When the Fault is triggered (i.e. when at least one of its associated trigger
groups is triggered),
the simulation variable is modified in some way.

\subsubsection{Fault Types}
\par The model includes 7 types of Faults:
\begin{itemize}
  \item Bias Faults
  \item Scale Faults
  \item Overwrite Faults
  \item Stale-value Faults
  \item White-noise Faults
  \item Random-walk Faults
  \item Function Faults \begin{itemize}
    \item Linear Faults
    \item Sine-Wave Faults
    \item Square-Wave Faults
    \item Triangle-Wave Faults
  \end{itemize}
\end{itemize}
See section \ref{sec:formulation:fault:fault_type} for descriptions of these
types.

\subsubsection{Supported Data-types}
Each of these 7 types has its own template class (the different types of
function faults are options within the same class).
They have all been tested on the following variable types:
\begin{itemize}
  \item Integers \begin{itemize}
    \item \verb"char" (treated as an 8-bit integer)
    \item \verb"unsigned char"
    \item \verb"short"
    \item \verb"unsigned short"
    \item \verb"int"
    \item \verb"unsigned int"
    \item \verb"long"
    \item \verb"unsigned long"
    \item \verb"long long"
    \item \verb"unsigned long long"
  \end{itemize}
  \item Floating-point Numbers \begin{itemize}
    \item \verb"float"
    \item \verb"double"
  \end{itemize}
  \item Boolean values (Overwrite and Stale-value faults only)
\end{itemize}
The Fault Manager only allows these types of variables (and references to them)
to be faulted; of particular note, pointers (to any data type) are not a
legitimate variable type.






\chapter {Requirements}

\section{Fault Effects}
\subsection {Profiles}
The model shall support the ability to modify a simulation variable with
either, or both:
\begin{enumerate}
 \item introduce a one-time initial Fault to modify the initial value of
       some variable, applied during model initialization.
 \item introduce a time-evolving profile that is consistent with commonly
       seen faults. These include, at a minimum:
 \begin{enumerate}
  \item Holding a variable to a pre-determined value
  \item Incrementing the nominal evolution of a variable with some fixed bias
  \item Scaling the nominal evolution of a variable by some fixed scale-factor
  \item Setting a value from a variable's current value, and holding the
        variable to that value.
  \item Introducing a noise signature over a variable's nominal evolution
  \item Introducing a walk-off signature over a variable's nominal evolution
  \item Introducing a random-walk accumulation over a variable's nominal
        evolution.
  \item Introducing a periodic-fluctuation signal over a variable's nominal
        evolution. See below for additional requirements on this capability.
 \end{enumerate}
\end{enumerate}
The model shall also provide an extensible framework to support user-defined
Fault signatures not currently required.

\subsection {Periodic Profiles}
 The last item of the list above introduced the requirement for periodic
 profiles, which have their own set of requirements.
 These periodic functions shall use some independent variable -- that is not
 necessarily associated with time -- to
 control the evolution of the periodic signature.
\begin{itemize}
 \item The signatures shall include, at a minimum:
  \begin{itemize}
   \item sine-wave
   \item triangular-wave, a linear approximation to a sine-wave
   \item square-wave, a discrete approximation to a sine-wave.
  \end{itemize}
 \item The independent variables shall be selectable from a variety of data
       types, including:
  \begin{itemize}
   \item integer of multiple forms (int, unsigned int, etc.)
   \item floating-point values of multiple forms (float, double, etc.)
   \item character-based variables (including individual character and string)
  \end{itemize}
 \item The parameters defining the periodic profile shall be configurable.
 \item The parameters defining the periodic profile shall be representable as
       their own independent functions.
\end{itemize}

\subsection{Compounding Faults}
The Fault signatures identified in the previous section shall be compoundable --
that is the model shall support combinations of these effects. For example, a
particular Fault might completely overwrite the evolution of some variable into
a noisy signature about some constant value.

\subsection {On and Off}
The model shall provide the ability to turn the Fault on and off repeatedly.
This could be either:
\begin{itemize}
 \item periodically, with the fault signature being applied for pre-determined
 durations at regular intervals during the simulation
 \item ad-hoc or quasi-randomly, with the fault signature being applied as
 needed
\end{itemize}

\subsection{Monte Carlo Support for Random Faults}
The model shall provide a mechanism that supports Monte Carlo based distribution
of the fault parameters. This is particularly important for dispersing the seed
when representing faults that naturally use random number generation, such as
introducing white-noise to a variable.

\subsection{Modifying Faults}
The model shall provide the capability to modify the signature of a Fault while
it is being applied. A simple example of this might be to change the value to
which a variable is being held.

\section{ Fault Causes and Triggers}
\subsection{Logical Operators}
The model shall support a wide set options for determining when a fault should
be applied. These tests shall support the following logical comparators (and
their near-inverses):
\begin{itemize}
 \item is equal to / is not equal to
 \item is greater than / is less than or equal to
 \item is less than / is greater than or equal to
\end{itemize}

\subsection{Logical Combinations}
The model shall provide control for triggering a Fault subject to multiple
constraints, with the ability to combine constraints with both AND and OR
logical operations, applied to single constraints and to groups. For example
\begin{itemize}
 \item A and B
 \item A or B
 \item (A and B) or (C and D)
\end{itemize}

\subsection{Triggering Variables}
\begin{enumerate}
 \item The model shall support a wide set of variable types to use as triggering
   mechanisms, including:
  \begin{itemize}
   \item boolean
   \item integer of multiple forms (int, unsigned int, etc.)
   \item floating-point values of multiple forms (float, double, etc.)
   \item character-based variables (including individual character and string)
  \end{itemize}
  Note -- this requirement is best achieved by making the trigger mechanism be
  template-based. While this is not a requirement per-se, it is a natural
  solution
  to the actual requirement.
 \item The model shall provide the capability to configure and later modify the
   value of the trigger-variable that results in the trigger-condition being
   satisfied.
\end{enumerate}

\subsection {Triggering Controls}
The model shall provide additional mechanisms for controlling recognition of a
successful trigger, including:
\begin{itemize}
 \item Enabling and disabling specific triggers (when a Fault has multiple
 triggers)
 \item Limiting the number of times a trigger can ``fire''
 \item Limiting the duration of an active trigger in terms of its independent
   variable advancing, making a trigger periodic such that it fires for a
   defined interval followed by an interval of not firing, repeated cyclically
   and indefinitely.
 \item Allowing a trigger to fire every time the trigger-condition is
   satisfied, for an indefinite period.
 \end{itemize}

\section {Fault Application}
The model shall support modifying any fundamental data type with an
appropriately configured Fault. Data types will include, at a minimum:
\begin{itemize}
 \item boolean
 \item integer of multiple forms (int, unsigned int, etc.)
 \item floating-point values of multiple forms (float, double, etc.)
 \item character-based variables (including individual character and string)
\end{itemize}

The model shall also provide an extensible framework to support application to
user-defined data-types.

\section {Fault Implementation and Configuration}
The previous requirements have addressed the necessary flexibility in
configuring the specific behavior of a Fault. This section addresses the
implementation of the Fault-management model within a simulation, and its
interface to both the user and the simulation. These are two separate
considerations:
\begin{itemize}
  \item how the faults are introduced into, and applied within,
  the simulation (i.e. the interface between the Fault instance and the model
  data it is affecting), and
  \item how the Fault instances are created and configured (i.e. the interface
  between the Fault instance and the simulation developer/user).
\end{itemize}

There are two distinct design patterns for how the faults may be implemented
in the simulation architecture, and two patterns for the
configuration-interface to those faults.

Because the implementation and configuration are separable aspects of
introducing faults, the Fault-management model shall support the use of either
implementation method with either configuration method.
We will look at these in turn -- the
implementation options and then the configuration options.

\subsection{Fault Implementation / Execution}
The Fault-management model shall support implementation of faults as:
\begin{itemize}
 \item an integrated component of a model, built into the
 design of the model, compiled with the rest of the model code, and executed
 in-line with the routine updates of that model.
 \begin{itemize}
  \item This requires that there be some
  mechanism to identify which Fault instances should be applied at which
  points within a model's routine update process to achieve the desired
  fault signature.
 \end{itemize}
 \item a component or set of components at the top-level of the simulation,
 instantiated independently of the instantiations of the models to which
 the faults will be applied. Execution of
 such a component will result in the application of its associated Fault
 instances to their associated variables.
 \begin{itemize}
  \item This requires that the Fault model have access to those variables to
  which the fault will be applied; these will be located elsewhere within
  the simulation.
 \end{itemize}
\end{itemize}

Note that these are mutually incompatible for a given fault and its
associated Fault instance, but not between faults (or Fault instances)
within a simulation. The model shall provide the flexibility to support
using both of these options within a single simulation.

The first of these implementation options provides a greater degree of
flexibility as to
the specific location of the application of a fault within the model's internal
algorithms. This is typically the option preferred by model developers. For
example, it might be found in sensor models where a fault may reside
in the input data, another in the processing of the input data, and another in
the output from the model.

The second of these implementation options provides a non-invasive capability
to support injection of faults in an otherwise-established simulation. This is
typically the option preferred by simulation managers and trainers, especially
where models have been constructed and compiled without an
integral fault mechanism. In these situations, it would be unreasonably onerous
to break the wrapper on a verified model to introduce a fault to it, so the
fault mananagement model must provide the capability to inject a fault into
the input, output, or some internal variable of the target model without being
invasive to the model execution. For example,
this implementation might be found at the interface between the simulation
environment and the flight-software environment, allowing faults to be applied
to sensor data flowing from the simulation to flight-software, and to the
responsive instructions flowing back to the simulation from flight-software.

\subsection {Fault Configuration}
The Fault-management model shall support configuration of faults via:
\begin{enumerate}
 \item Direct access to separately instantiated Fault instances.
 \item Indirect access via a fault-manager.
\end{enumerate}

Notes on interpretation:
\begin{itemize}
 \item unlike the implementation options, these are \textbf{not} conceptually
   incompatible for a given fault.
   CML contains several examples of models (unrelated to this Fault-management
   model) that include a manager-class to manage separately instantiated
   entities, each of which gets registered with the manager. The same
   architecture could be implemented here, to have individually instantiated
   Fault instances be registered with a fault-manager, thereby allowing access
   through the manager and to the individual instantiations directly.
 \item However, that instantiate-and-register operational paradigm is not
   specifically called out in these requirements, and has not been implemented
   in the Fault management model at this time.
 \item The model does provide a fault-manager, which does maintain a
   list (STL-\textit{list}) of the faults for which it is responsible. This
   fault-manager is responsible for creating those Fault instances for which
   it is responsible. The model \textit{does
   not} provide a mechanism for registering separately instantiated Fault
   instances with a manager.
 \item The implemented interpretation of this requirement ties together the
   instantiation and management, i.e. a Fault shall be either
   instantiated and managed individually, or instantiated and managed by the
   manager.
 \item Within a single simulation, there may simultaneously be a Fault-manager
   and individual Fault instances independent of that manager. There may be
   multiple managers, each responsible for multiple dependent Fault instances,
   and multiple independent Fault instances.
 \item This point may be revisited at some point in the future.
\end{itemize}

\chapter {Formulation}
\section{Definitions}
\begin{itemize}
  \item \textit{\textbf{triggered}} This is an overloaded term:
  \begin{itemize}
    \item A \textit{Trigger} is described as triggered when its logical
      conditions are satisfied.
    \item A \textit{Trigger-group} is described as triggered when all of its
      associated \textit{enabled} Trigger instances are individually triggered.
      Any Trigger that has been disabled in a Trigger-group is not included
      in th evaluation of whether that Trigger-group is triggered.
    \item A \textit{Fault} is described as triggered when any of its associated
      Trigger-group instances are triggered AND the Fault is enabled.
    \end{itemize}

  \item \textit{\textbf{reset}} A Fault is reset whenever it transitions from
    being not-triggered to being triggered. This is typically a consequence of
    the transition of one of its Trigger-groups, but could also as a consequence
    of enabling the Fault while the Trigger-group conditionals are satisfied.
  \item \textit{\textbf{target-variable}} This refers to the variable -- found
    elsewhere in the simulation -- to which we want to apply a fault. The value
    of this variable will be modified by the application of the fault.
  \item \textit{\textbf{modifier}} When the fault being applied is additive,
    i.e. it adds some value to the nominal value of the target-variable, the
    modifier is the value that is added to the nominal value of the
    target-variable.
  \item \textit{\textbf{trigger-variable}} This refers to the variable -- found
    elsewhere in the simulation -- that will be monitored for determining when
    the fault is applied.
  \item \textit{\textbf{trigger-value}} This is the value of the
    trigger-variable that results in the application of the fault.
    The Fault instance representing the desired fault may be
    configured to trigger when the value of the trigger-variable is related in
    some specified way (e.g. equal to, or less than, or greater than, etc.) the
    trigger-value.
  \item \textit{\textbf{firing / activating}} In places in this document, the
    act of having a Fault modify the target variable is referred to as the
    Fault firing or activating.
  \item \textit{\textbf{fire-count}} This is a count of how many times some
    Fault has fired or of how many times some Trigger has triggered.
\end{itemize}

\section {Fault Supporting-concepts}

\subsection{Limiting Number of Fault Applications}
\label{sec:formulation:fault_support:fire_count}
Faults have an option to limit the number of times they ``fire'', i.e.
apply a modification to the target variable.
If this option is enabled, the Fault counts how many times it applies a
modification and disables itself when the count reaches a certain limit. Note:
\begin{itemize}
 \item This is a count limit, not a time period; if the intent is to apply
   the Fault for some period of time, a better strategy would be to add a
   time-dependent trigger and close out the fault through the trigger mechanism;
   using the count brings in an undesirable dependency on the rate at
   which the Fault instance is updated; unexpected consequences may be
   observed if the update-rate is changed.
 \item The count is zeroed whenever the Fault is reset. The count can be zeroed
   any number of times.
 \item After being disabled by reaching its count limit, the Fault can be
   re-enabled manually. At that time, the Fault's \textit{triggered} status will
   be re-evaluated and if it is still triggered, the Fault instance will be
   reset and the
   count zeroed (see previous bullet). Until the Fault is manually re-enabled,
   the conditional tests associated with the triggers will not be evaluated.
\end{itemize}

\subsection{Limiting Number of Trigger Applications}
\label{sec:formulation:fault_support:trigger_fire_count}
Like a Fault, a Trigger can also be limited to the number of consecutive
times it can be considered ``triggered''.
If this option is enabled, the Trigger counts how many times it has been
used with a positive status -- i.e. the number of times it has been tested
against its threshold and found to meet the criterion for triggering.
When the count reaches a user-specified limit, the Trigger reverts to a
non-triggered state regardless of whether the triggering criterion continues
to be satisfied.

Note that:
\begin{itemize}
 \item This is a count limit, not a time period; if the intent is to apply
   the Trigger for some period of time, a better strategy would be to add a
   second trigger to the same group, with the criterion that time be less
   than some value. Remember that all Triggers within a group must be
   triggered for the group to be triggered.
   Using the count brings in an undesirable dependency on the rate at
   which the Trigger's conditional criterion is evaluated, and on the number
   of Faults using the Trigger;
   unexpected consequences may be observed if the update-rate is changed.
 \item The count is zeroed if the Trigger returns a negative state for any
   reason \textit{other than} the count exceeding the limit.
 \item A Trigger blocked by an excessive count is still considered enabled;
   it will block the \textit{TriggerGroup} from returning a positive state
   and it will continue to be evaluated on every cycle.
 \item Remember that a single Trigger can be used in multiple Faults.
   Limiting a Trigger will potentially affect all of the Faults which use that
   Trigger. Additionally, the Trigger will be tested by each Fault,
   incrementing the use-count with each test. If the count-limit is not an
   integer multiple of the number of Fault instances using the Trigger, at
   some point there will be a cycle with some subset of the Fault instances
   using this Trigger receive a positive response from the Trigger, while
   others get a negative response after the Trigger has been count-limited.
\end{itemize}

\subsection{Fault Functions} \label{sec:formulation:concept:functionss}
For some faults, the value of the modifier is a function of one or more
independent variables, e.g. $y=mx+b$, $y=A sin(\omega t + \phi)$
See Section \ref{sec:formulation:faults:function} for more details on the
operation and configuration of Fault Functions. For now, we focus on the
concepts.

All Fault Functions are represented as an instance of a templated-type,
\verb'FaultFunction', which inherits from the non-templated
\verb'FaultFunctionBase', which in turn inherits from two classes:
\begin{itemize}
\item \textbf{Fault} A FaultFunction is a type of Fault, so inherits all the
capabilities of a \verb'Fault'.
\item \textbf{FaultFunctionParameter} requires a little more explanation.
The parameters in some fault-functions (e.g. the
amplitude $A$ in $y=A sin(\omega t + \phi)$) may themselves be linear
functions of some (other) independent variable. In this case, the label
\textit{``parameter''} is something of a misnomer; \textit{``variable''}
would be a more accurate description, with the value of the
\textit{``parameter''} being a
linear function ($y=mx+b$) of two parameters (
\verb'rate', $m$ ;  \verb'nominal', $b$) and an independent variable
(\verb'ind_variable', $x$). We create a class \verb'FaultFunctionParameter'
to provide this linear-function capability.

Note, that the Fault Function itself must also have this linear
function capability to represent linear functions so is, in effect, an
extension of a \verb'FaultFunctionParameter'. Thus, we can use the
\verb'FaultFunctionParameter' class as an integral component of the
\verb'FaultFunction' class by making \verb'FaultFunctionParameter' be
inherited into the \verb'FaultFunctionBase' class itself. The
\verb'FaultFunction' class-template then uses the components of its inherited
\verb'FaultFunctionParameter' to represent the core components of its
function-nature, while using additional instances of
\verb'FaultFunctionParameter' to
represent the values of the parameters/variables associated with its function.
\end{itemize}
This gives rise to two more concepts, intrinsic to understanding fault-function
-- that of the fault-function parameter, and the fault-function independent
variable.



\subsubsection{Fault-Function Parameters} \label{sec:formulation:concept:params}
The parameters used for evaluating fault-functions are
represented by the class \verb'FaultFunctionParameter'.

This class includes the capability to evaluate a linear expression $y=mx+b$. In
this expression, $m$, $x$, and $b$ are class members.
\begin{itemize}
 \item \verb'nominal' $b$
 \item \verb'rate' $m$
 \item \verb'ind_variable' $x$; this is an instance of
   \verb'FaultFunctionIndependentVariable', described in the next section.
\end{itemize}

Note -- remember we are specifically discussing the
\verb'FaultFunctionParameter'
aspect of a fault-function, not the entire structure of the fault-function.
Fault-functions are not limited to any specific type of function, but the value
of its parameters is restricted to either being constant (i.e. a true parameter)
or a linear expression. See section \ref{sec:formulation:faults:function} for
more details on the operation and configuration of Fault Functions.

More complicated fault-functions use other parameters; each of these can be a
linear-function of an independent variable, not necessarily the same
independent variable as used by the fault. Thus, each of these
\textit{``parameters''}
may be represented by an instance of \verb'FaultFunctionParameter'.

For example, a fault might generate a modifier $y$ with a sine-function
dependency on $x$:
\begin{equation*}
 y= Asin(\omega x + \phi)
\end{equation*}

The \textit{parameters} in this expression -- $A$, $\omega$, and $\phi$ --
could each be an independent linear function of some other variable.
Expanding each of these parameters into its own linear function
yields:
\begin{equation*}
y= (m_A x_A + b_A) sin((m_\omega x_\omega + b_\omega)x +
(m_\phi x_\phi + b_\phi))
\end{equation*}

This concept is explored in more detail, including the process for setting
these parameters, in Section \ref{sec:formulation:fault:params}.

\subsubsection{Independent Variables} \label{sec:formulation:indep_var}
Independent variables are used in Fault-functions (see Section
\ref{sec:formulation:faults:function}) to represent the independent-variable
used in evaluating the modifier. In code, these are represented by instances of
the class \verb"FaultFunctionIndependentVariable".
Each instance of \verb"FaultFunctionIndependentVariable" accesses the value of a
numeric
variable residing elsewhere in the simulation; the nature of this access
requires some explanation:
\begin{itemize}
  \item To support on-the-fly creation and population via XML,
    \verb"FaultFunctionIndependentVariable" instances cannot be constructed with
    a direct
    reference to the independent-variable, the accessor must be a pointer,
    assigned in a post-construction \textit{initialize(...)} method.
  \item To support multiple data types, the pointer cannot be defined as
    type-specific.
  \item To support multiple data types \textit{and} delay the population of the
    pointer to initialization, the pointer cannot even by typed from a template
    (which would result in the type being specified at construction).
  \item Instead, the \verb"FaultFunctionIndependentVariable" class uses the
    \verb"UntypedVariableBase" class, and the template-class
    \verb"UntypedVariable" that inherits from \verb'UntypedVariableBase'.
    \begin{itemize}
      \item The \verb"FaultFunctionIndependentVariable" class has a pointer
        (called \verb'variable') of type \newline
        \verb'UntypedVariableBase *'
        that addresses some instance of
        \verb'UntypedVariable'.
      \item At initialization of the \verb"FaultFunctionIndependentVariable",
        the argument
        passed in is used as a template
        parameter to create a \textbf{new} instance of \verb'UntypedVariable',
        and the \verb'variable' pointer is assigned to that new
        instance.
      \item \verb'UntypedVariable' contains a reference (also called
        \verb'variable') to the desired independent-variable found elsewhere in
        the simulation. This independent-variable is passed as the argument to
        \verb'FaultFunctionIndependentVariable::initialize(T &)') and from there
        is passed as the argument to the constructor of the new
        \verb'UntypedVariable' instance. \verb'UntypedVariable' can support a
        reference to the
        independent-variable because the details (type and address) of the
        independent-variable are both known at the construction of the new
        \verb'UntypedVariable'.
      \item the memory-address of the new instance of
        \verb'UntypedVariable' is not particularly important, and it certainly
        does not have to be contiguous with the fault. We do need to maintain
        access to this new instance; its address is recorded as the
        \verb"variable" pointer found
        in \verb"FaultFunctionIndependentVariable".
    \end{itemize}
    \item To obtain the value of the independent-variable, we use:
      \begin{itemize}
        \item \verb'(FaultFunctionIndependentVariable::get_value()' which calls
        \item \verb'(FaultFunctionIndependentVariable::variable)->get_value()',
          which accesses the appropriate instance of
          \verb'UntypedVariable::get_value()'
        \item \verb'UntypedVariable::get_value()' static-casts the value
          of its \verb'variable' reference (to the actual independent-variable)
          from its native type (determined by template parameter) to type
          \verb'double', and returns that value to \newline
          \verb'FaultFunctionIndependentVariable::get_value()'
        \item \verb'FaultFunctionIndependentVariable::get_value()' branches
          depending on
          whether the variable is being treated as absolute or relative (see
          discussion immediately following) and returns the resulting value to
          the calling function.
      \end{itemize}
    \item Caveat: the template parameter used in the template-class
      \verb'UntypedVariable' must support a static-cast to type \verb'double'.
\end{itemize}

A \verb"FaultFunctionIndependentVariable" instance can be configured to
represent the value
of the independent-variable as either its true value, or as a value relative to
its initial value.
\begin{itemize}
  \item An \textit{absolute} configuration (\textit{relative\_value = false})
    simply provides the current value of the independent-variable back to the
    calling function.
  \item A \textit{relative} configuration (\textit{relative\_value = true},
    default) requires a record of the value of the independent-variable when
    the fault was most recently reset. The value provided to the calling
    function is then the difference between the current and recorded values
    of the independent-variable.
\end{itemize}




\section {Fault Types} \label{sec:formulation:fault:fault_type}
\subsection {Nomenclature}
\begin{itemize}
  \item $x$ represents the variable to which the fault is applied
  \begin{itemize}
    \item $x_i$ is the value of the variable at the point of interest
    \item $x_0$ is the value of the variable at the point at which the fault was
      most recently reset (both triggered and enabled).
    \item $\Delta x$ is the modifier, the value added to the variable when
      applying an additive fault.
    \item $x'_i$ is the modified value of $x_i$; in some cases, $x'_i = x_i +
      \Delta x$.
  \end{itemize}
\end{itemize}

\subsection{Bias Faults}
\par When a bias fault is triggered, it adds a constant bias \textit{b} to the
variable.
\begin{equation}
x'_i = x_i + b
\end{equation}

\subsection{Scale Faults}
\par When a scale fault is triggered, it multiplies the variable by a scale
value \textit{s}.
\begin{equation}
x'_i = s x_i
\end{equation}

\subsection{Overwrite Faults}
\par When an overwrite fault is triggered, it overwrites the variable with some
pre-determined value $x_f$.
\begin{equation}
x'_i = x_f
\end{equation}

\subsection{Stale-value Faults}
\par When a stale-value fault is triggered, it holds the variable at its current
value until it is no longer triggered.
\begin{equation}
x'_i = x_0
\end{equation}

\subsection{White-noise Faults}
\par When a white-noise fault is triggered, it adds white noise $\eta$ to the
variable. The noise may have a Gaussian or uniform distribution.
\begin{equation}
x'_i = x_i + \eta_i
\end{equation}

\subsection{Random-walk Faults} \label{sec:intro:randomwalk}
\par When a random-walk fault is triggered, it adds random-walk error to the
variable. Essentially, the fault contains a bias value \textit{w} that is
initialized to 0 every time the fault is reset and is incremented by a random
value $\eta$ every timestep the fault remains triggered.
The random value $\eta$ may be configured to have a Gaussian or Uniform
distribution.
If the fault status changes to non-triggered, the bias value is reset to 0.
\begin{equation*}
w_i = w_{i-1} + \eta_i
\end{equation*}
\begin{equation}
x'_i = x_i + w_i
\end{equation}

\subsection{Function Faults} \label{sec:formulation:faults:function}

With Function Faults, the modifier $\Delta x$ is computed as some function of a
second, independent variable. In this section, the independent variable is
represented with the symbol $t$ because the most commonly used independent
variable is time. However, it is important to comment that $t$ in this context
does not represent time exclusively, it is simply convenient and could represent
any variable.
\begin{itemize}
  \item $t_i$ is the value of the independent variable at the point of
    interest
  \item $t_0$ is the value of the variable at the point at which the fault was
    most recently reset (both triggered and enabled).
\end{itemize}

\subsubsection{Linear Function Faults}
\par Linear function faults use two parameters called \textit{rate, m} and
\textit{nominal, b}.
When the fault is triggered, the fault variable is biased by a linear function
of the independent variable.
\begin{equation*}
\Delta x_i = m t_i + b
\end{equation*}
\begin{equation} \label{eqn:linfunc}
x'_i = x_i + m t_i + b
\end{equation}

\subsubsection{Periodic Function Faults}
\par Periodic function faults represent the modifier $\Delta x$ as a periodic
function $g(\phi)$, i.e.:
\begin{equation*}
  \Delta x = A g(\phi)
\end{equation*}
where:
\begin{itemize}
 \item $g(\phi)$ is a normalized periodic function, with amplitude and period
   both 1.0; i.e. $g(\phi) \in [-1.0, 1.0]; g(n + \phi) = g(\phi)$ with
   $n \in \mathbb{Z}$.
 \item $A$ represents the amplitude of the modifier and may itself be
   variable, in which case $A$ would be better represented as $A_i$, the
   amplitude at step $i$.
 \item $\phi$ represents the phase of the function, and depends on the
   independent variable and two additional variables:
   \begin{itemize}
     \item the function frequency, $f$ (which may also be variable, a linear
       function of some other independent variable); this has
       dimension that is inverse to the independent variable (e.g. with the
       independent variable representing time with units of seconds, $f$ would
       have units of Hz).
     \item the phase offset $\psi$. This is typically constant with value
       equal to $\phi_0$, the value of $\phi$ at the time the fault was most
       recently reset. However, it could be (unusually) an independent
       linear function of some other variable.
    \end{itemize}
    Note -- the phase and phase-offset in this model are based on
    whole-cycles rather than a radian-based representation.
    Cycles repeat with a phase difference
    of $1.0$, \textbf{not} with a phase difference of $2 \pi$.
\end{itemize}


The value of $\phi$ is computed by integrating the frequency
\begin{equation*}
  \phi_i = \psi_i + \int_{t_0}^{t_i} f \mathrm{d}t
\end{equation*}
Then the modifier and final value for step $i$ are:
\begin{align*}
  \Delta x_i &= A_i g(\phi_i) \\
  x'_i &= x_i + \Delta x_i
\end{align*}

As a simple example, consider $f$, $A$, and $\psi$ as constant, with $\psi =
\phi_0 = f t_0$.
Then we get the simple form:
\begin{equation*}
  x'_i = x_i + A g(f t_i)
\end{equation*}

\subsubsection{Periodicity Simplification}
Because $g(\phi)$ is periodic, we only need to evaluate $g$ in the unit domain,
$[-0.5, 0.5)$. For simplicity, we will work with
$\phi' \in [-0.5,0.5)$,
the difference betweeen $\phi$ and the closest integer. Algorithmically,
\begin {equation*}
  \phi' = \phi - round(\phi)
\end {equation*}
Note: we can write $g(\phi)$ as a function of $\phi'$, where $\phi$ is the
unconstrained phase angle and $\phi'$ is its corresponding constrained
value.

Currently, the model supports three periodic functions:
\begin{itemize}
  \item \textbf{sine wave:}
    \begin{equation*}
      g(\phi) = sin(2 \pi \phi')
    \end{equation*}

 \item \textbf{square wave} is a discrete approximation to the sine wave,
    taking a value either $\pm 1$.
    \begin{equation*}
      g(\phi) = \left\{ \begin{array}{ll}
                 -1 & : \phi' < 0 \\
                 +1 & : \phi' \geq 0 \\
                \end{array} \right.
    \end{equation*}

  \item \textbf{triangle wave} is a linear approximation to the sine-wave,
    \begin{equation*}
      g(\phi) = \left\{ \begin{array}{ll}
                       -4 \phi' - 2  & : \phi' < -0.25 \ \\
                        4 \phi'      & : -0.25 \leq \phi' < 0.25 \\
                       -4 \phi' + 2  & : \phi' \geq 0.25
                     \end{array} \right.
    \end{equation*}
    This may also be thought of as a modulated absolute-value function, phase
    shifted by 0.25 to the left before applying the $round$ function,
    yielding the simpler form:
    \begin{equation*}
      g(\phi) = 4 |(\phi +0.25) - round( \phi+0.25) | -1
    \end{equation*}
\end{itemize}

The periodic parameters (frequency, amplitude, and phase offset) can each be
specified as a linear function of some independent variable.
Their value is computed from the standard $y=mx+b$, with $x$ in this context
represented by an instance of \textit{FaultFunctionIndependentVariable} (see
section \ref{sec:formulation:indep_var}) and can be configured to represent
either the value of the independent variable, or its value relative to its
initial value. $m$ and $b$ in this context are identified as \textit{rate} and
\textit{nominal} respectively and can be set with the \textit{set\_param}
method (see Section \ref{sec:formulation:fault:params}).
\begin{align*}
  f &= f_0 + m_f t_f \\
  A &= A_0 + m_A t_A \\
  \psi &= {\psi_0} + m_{\psi} t_{\psi}
\end{align*}

If $f$ is constant, the phase
\begin{equation*}
  \phi_i = \psi_i + \int_{t_0}^{t_i} f \mathrm{d}t  \;=\;
           \psi_i + f_0(t_i - t_0)
\end{equation*}

However, if $f$ is variable ($f=f_0 + m_f t_f$), the phase
\begin{equation*}
  \phi_i = \psi_i + \int_{t_0}^{t_i} f \mathrm{d}t  \;=\;
           \psi_i + f_0(t_i - t_0) + m_f \int_{t_0}^{t_i} t_f \mathrm{d}t
\end{equation*}
which requires knowledge of the relationship between $t_f$ and $t$ (which are
not necessarily the same variable). This is typically not available so the model
uses a linear approximation between each timestep:
\begin{equation} \label{eqn:tfinterp}
t_f (t) \approx {t_f}_{i-1} + \frac{{t_f}_i - {t_f}_{i-1}}
                                    {t_i - t_{i-1}}
                                    (t - t_{i-1})
        \;:\; t_{i-1} \leq t < t_i
\end{equation}
which allows the integral to be evaluated in pieces using trapezoids, with each
trapezoid contributing:
\begin{equation*}
\int_{t_{i-1}}^{t_i} t_f \, \mathrm{d}t \approx \frac{{t_f}_i + {t_f}_{i-1}}{2}
                                                (t_i - t_{i-1})
\end{equation*}
and consequently
\begin{equation*}
 \int_{t_{i-1}}^{t_i} f \, \mathrm{d}t
    \approx \frac{f_i + f_{i-1}}{2}  (t_i - t_{i-1})
\end{equation*}

The numerically-integrated frequency can then be accumulated:
\begin{align*}
  F_i &= \int_{t_0}^{t_i} f \, \mathrm{d}t \\
    &=  F_{i-1} + \int_{t_{i-1}}^{t_i} f \, \mathrm{d}t \\
    &\approx F_{i-1} + \frac{f_i + f_{i-1}}{2}  (t_i - t_{i-1})
\end{align*}
The phase can then be expressed in terms of $F$:
\begin{equation*}
  \phi_i \approx \psi_i + F_{i-1} + \frac{f_i + f_{i-1}}{2}  (t_i - t_{i-1})
\end{equation*}

Indeed, to avoid excessive logical branching to check on the value of $m_f$
(\verb'frequency.rate') at each cycle, this formulation is used in all cases,
including those for which $f$ is constant (i.e. $m_f = 0$).

\subsubsection {Negative Periodic Parameters}
The computation of these parameters as linear functions introduces the potential
for the parameters to take negative values.\textbf{No checks are made on the
sign of these parameters}. The rationale to support this design decision is
that, while we conventionally consider amplitude and frequency to be positive
values, the formulation presented above does not preclude negative values.
\begin{itemize}
 \item \textit{phase\_offset} Can be permitted to be negative without
   complication. Phase-offsets are periodic, with $\psi$ being equivalent to
   $n+\psi, \forall n \in \mathbb{Z}$.
 \item \textit{frequency} can be permitted to be negative because we use the
   frequency to
   advance the phase; a negative frequency combined with an increasing
   independent
   variable simply results in a decreasing phase. Note that this is equivalent
   to having a positive frequency with a decreasing independent variable.
 \item \textit{amplitude} can be permitted to be negative because a negative
   amplitude simply means that the modifier has a different sign than it would
   with a positive amplitude of the same magnitude. It could be argued that
   changing the sign of the
   amplitude would introduce a sign change in the value of the modifier which
   could lead to an unexpected discontinuity. However, with the amplitude being
   a linear function, it can only change sign by passing through 0, at which
   point the value of the modifier goes to zero as well.
\end{itemize}


\subsection{Function Fault Parameters} \label{sec:formulation:fault:params}

The parameters used for evaluating function-faults -- namely the nominal
and rate values for a linear-function-fault, and the frequency, amplitude, and
phase-offset for a periodic function-fault -- can be set directly, or via the
\textit{FaultFunctionBase::set\_param(...)} method. This method takes two
arguments, with an optional third argument that is explained below. The
arguments are:

\begin{itemize}
 \item A character string, describing which parameter is to be set. These can be
  \begin{itemize}
   \item initial or nominal (these are equivalent)
   \item rate
   \item frequency
   \item frequency\_rate
   \item amplitude
   \item amplitude\_rate
   \item phase\_offset
   \item phase\_offset\_rate
  \end{itemize}
 \item The new value for the specified variable.
 \item (optional) A boolean flag indicating whether a new assignment to a rate
   variable should result in an automatic adjustment to its corresponding
   nominal value. This flag defaults to false if not provided.
\end{itemize}

In the case of setting the \textit{nominal} parameter, or setting the
\textit{rate} parameter without the optional flag, the new value is simply
pushed onto the variable. The optional flag is provided as a convenience when
adjusting the rate mid-fault.

The resulting value of the linear function -- whether it be the value of the
modifier itself when the fault is configured to follow a linear function, or the
value of one of the periodic parameters which also follow a linear function --
is computed from $y = mx +b$, with $m$ being defined by \textit{rate}, $b$ being
defined by \textit{nominal}, and $x$ being the associated independent variable.
(see Section \ref{sec:formulation:concept:params}).
If we change the rate $m$, the value $y$ will change accordingly unless the
value $b$ is also changed; this would produce a discontinuity in the value (see
the red line in Figure \ref{fig:Rate_change}).

When changing the rate $m$, avoiding this discontinuity necessitates changing
the nominal value $b$. However, this is not a universally
desirable outcome, so the capability is provided to adjust $m$ with and without
a corresponding change to $b$. The choice is identified by setting the optional
flag.

We have the output value $y$ (before modification) and $y'$ (after modification)
written as
\begin {align*}
 y &= mx + b \\
 y' &= m'x + b'
\end {align*}

If we want the modifier to be continuous, we require:
\begin {align*}
 y' &= y \\
 b' &= y - m'x
\end {align*}
The consequences of adjusting the rate with and without the optional flag are
shown in Figure \ref{fig:Rate_change}.

\begin{figure}
 \begin{center}
 \includegraphics[scale=0.5]{graphics/Rate_change.png}
 \caption{ The three dashed lines show reference linear functions:\newline
    $y = x+100, y=-x+100, y=-x+200$.\newline
    Consider a fault starting at $x=0$ with \textit{nominal} value $100$ and
    \textit{rate} $1$. It is applied continually with these parameters until
    $x=50$.
    At $x=50$ the rate is changed from $+1$ to $-1$.\newline
    The red line illustrates
    the time evolution of this fault if the rate-change is made while leaving
    the \textit{nominal} value unchanged (default).\newline
    The black line illustrates the the time evolution of this fault if the
    rate-change is made with the optional flag set, which results in
    changing \textit{nominal} to $200$ in this case.}
  \label{fig:Rate_change}
  \end{center}
\end{figure}


If setting the parameters directly, the automatic adjustment to $b'$ is not
available, that only happens when using
\textit{FaultFunctionBase::set\_param(...)} with the 3rd argument \textit{true}.
If using the Fault Manager architecture (see section
\ref{sec:UG:FaultManager}), the method
\textit{FaultManager::set\_fault\_param(...)}
requires a preceding argument specifying the name of the fault, followed by the
same 2 or 3 arguments specified above; the first argument identifies the fault
and the remaining arguments are passed directly through to the corresponding
\textit{FaultFunctionBase::set\_param(...)}.

\section{Triggers} \label{intro_trigger}
\par A trigger consists of a reference to a simulation variable; a value; and a
comparison operator, which can be one of the following:
\begin{center}\begin{tabular}{|l|c|c|} \hline
\textbf{Operator} & \textbf{Symbol} & \textbf{TriggerBase Name} \\ \hline
Less than                & <      & LT \\ \hline
Less than or equal to    & $\leq$ & LE \\ \hline
Equal to                 & =      & EQ \\ \hline
Greater than or equal to & $\geq$ & GE \\ \hline
Greater than             & >      & GT \\ \hline
Not equal to             & $\neq$ & NE \\ \hline
\end{tabular}\end{center}
The trigger is triggered when the condition "(variable) (operator) (value)" is
true.
For example, if the variable is \textit{obj1.var}, operator is \textit{LT} and
the value is $2$, the trigger is triggered
when $obj1.var < 2$.

\par Each fault contains pointers to at least one trigger group.
Each trigger group contains pointers to at least one trigger.
Each trigger in a trigger group can be enabled or disabled at any time (this is
done at the trigger group level; if a trigger is used by multiple groups,
disabling it in one group will not disable it in the others).
The term "triggered" is overloaded here: a trigger group is triggered only
when \textbf{all} of its \textbf{enabled} triggers are triggered; a fault is
triggered when \textbf{any} of its trigger groups are triggered.

\subsection{Variable Types}
\par Triggers are implemented as a template class, \verb"Trigger".
In theory, they should work for any type for which the <, <=, ==, >=, >, and !=
operators are valid; however, this may not always have the expected behaviour.
For example, if the trigger variable is a \verb"float*" (pointer to float), the
value being compared is the address of a float, not the value stored at that
address.

\par The template has been specialized for strings (the type
\verb"std::string").
For string triggers, only the EQ and NE operators are valid.
If the operator is EQ, the trigger is triggered when the variable string matches
the value string exactly; if the operator is NE, the trigger is triggered when
the strings don't match exactly.

\par \verb"Trigger" has been tested on the following types:
\begin{itemize}
  \item Integers \begin{itemize}
    \item \verb"char" (treated as an 8-bit integer)
    \item \verb"unsigned char"
    \item \verb"short"
    \item \verb"unsigned short"
    \item \verb"int"
    \item \verb"unsigned int"
    \item \verb"long"
    \item \verb"unsigned long"
    \item \verb"long long"
    \item \verb"unsigned long long"
  \end{itemize}
  \item Floating-point Numbers \begin{itemize}
    \item \verb"float"
    \item \verb"double"
  \end{itemize}
  \item \verb"bool" (only makes sense if the operator is EQ or NE)
  \item \verb"std::string"
\end{itemize}
The fault manager only allows these types of variables (and references to them)
to be used as trigger variables; pointers are not allowed.

\subsection{Periodic Triggers} \label{intro_periodic}

The intent of periodic triggers is to make the fault appear at regular, periodic
intervals, where
period in this context is some specified difference between two values of the
independent-variable. This is most commonly used when the independent variable
is a measure of time, and period then is interpreted in its conventional sense
-- as a specified separation in time -- but the same concept can equally apply
to other variables such as distance (in which case ``period'' in this context
would be better aligned with a conventional use of ``wavelength''), etc.

Note that Periodic triggers can theoretically be used on most, but not all
variable types. Of particular note, in the current implementation, triggers that
use a STL-string as its independent variable do not support periodic triggers.
The simple rationale is that it is difficult to conceptualize (unambiguously)
how to define periodic behavior as a string (e.g. trigger when this string is
``abc'', continue to trigger for ``def'', then every ``ghi'' thereafter).


The parameters for a periodic trigger are the period and the trigger length,
which both have the same units as the independent variable:
every \textit{period} units, the fault is triggered for \textit{length}
units. For example, every 5 seconds the fault is triggered for 2 seconds.

\par In mathematical terms, making a trigger periodic adds an extra condition
that must be met for the trigger to be considered \textit{triggered}.
The first time the regular condition is met, the value of the trigger variable
$s$ is saved as $s_0$; thereafter, the trigger will be considered
\textit{triggered} when the regular condition is satisfied AND:
\begin{equation*}
  | (s- s_0) \% p | \leq l
\end{equation*}
where $p$ and $l$ are the period and length respectively.
Note that the \verb"std::fmod" implementation ($\%$) returns a negative value
when $s-s_0 < 0$, hence the need for taking the absolute value of the result.
The trigger will be applied in the region of the period that is closest to
$s_0$, this is visualized in the Figure \ref{fig:periodic_trigger}.

\begin{figure*}[h!]
  \includegraphics[scale=0.72]{graphics/periodic_trigger.png}
  \caption{Conceptual triggering logic, with green representing conditions
  satisfied, and red not-satisfied. The top line shows a hypothetical condition
  that would be the result without periodic triggering; the middle line shows
  the periodic triggering starting at $s_0$ with period $p=5$ and length
  $l=2$; the
  lower line shows the combined result. The upper plot is for an increasing
  independent variable, and the lower plot for a decreasing independent
  variable.}
\end{figure*}\label{fig:periodic_trigger}

\chapter{User's Guide} \label{sec:UG}
\par Faults can be implemented in two ways: by coding them directly into the
simulation or by implementing a fault manager and writing faults in an XML file.
Generally, the XML approach is preferable, as it makes it very easy to add and
configure faults; however, the direct approach makes it easier to see what is
happening under the hood and may be simpler for adding just one or two faults.

This User's Guide is presented in 3 sections:
\begin{enumerate}
 \item Overview of concepts
 \item Implementation of concepts using the direct-implementation architecture
 \item Implementation of concepts using the fault-manager and XML architecture
\end{enumerate}


\section {Overview of Concepts}
\subsection{Components (Required and Optional)}
\subsubsection{Location of Implementation}
The model is usually implemented as a part of a model representing some sensor
hardware. For the direct instantiation architecture, each piece is instantiated
independently. For the fault-manager architecture, the pieces are created
on-the-fly according to the provided specification.

\subsubsection{ Trigger and Trigger Group}
The model requires that each fault have at least one trigger group, and that
each trigger group have at least one trigger. Details on the conceptual
development of triggers and trigger-groups can be found in Section
\ref{intro_trigger}.

Trigger groups are added to a fault using the \textit{Fault::add\_trigger
\_group(TriggerGroup\&)} method.
Triggers are added to a trigger-group using the
\textit{TriggerGroup::add\_trigger(TriggerBase\&)} method.
If using the Fault Manager architecture, these methods are called automatically.

An important concept, worth repeating, is that the relationship between
faults and
trigger-groups and between trigger-groups and triggers is many-to-many:
a trigger can be added to multiple trigger groups, and a trigger group can be
added to multiple faults; a fault can have multiple associated trigger-groups,
and a trigger-group can have multiple associated triggers. A fault is triggered
when at least one of its trigger groups is triggered; a
trigger group is triggered when all of its triggers are triggered.
All triggers must be given a variable, an operator, and a value.

\subsubsection {Faults} \label{sec:UG:components:faults}
\par For more details on any of the features described here, see Sections
\ref{sec:intro:faults} and \ref{sec:formulation:fault:fault_type}.

The Fault Management supports 7 different types of faults.
Each of these types is represented by a different template class:
\begin{center}\begin{tabular}{|l|l|l|} \hline
\textbf{Fault Type} & \textbf{Class} & \textbf{XML Type Name} \\ \hline
Bias & FaultBias & BIAS \\ \hline
Scale & FaultScale & SCALE \\ \hline
Overwrite & FaultOverwrite & OVERWRITE \\ \hline
Stale Value & FaultStale & STALE \\ \hline
White Noise & FaultWhiteNoise & WHITENOISE \\ \hline
Random Walk & FaultRandomWalk & RANDOMWALK \\ \hline
Function & FaultFunction & FUNCTION \\ \hline
\end{tabular}\end{center}
Regardless of type, all faults must have a variable and may have a name (the
name is required if using a fault manager).

\subsubsection {Fault Parameters}
Each type of fault (except stale-value) has additional parameters specific to
that type.
Most of them can be modified in input files through the fault's
\textit{set\_param} method, whose arguments are the name of the parameter (as a
string) and the value to set the parameter to.
The fault manager has a method called \textit{set\_fault\_param} that acts as a
pass-through to \textit{set\_param}; its arguments are the name of the fault
followed by the \textit{set\_param} arguments. See section
\ref{sec:formulation:fault:params} for discussion of the consequences of setting
parameters.



\subsubsection{Independent Variables}
Independent variables are used in function-faults; these must be associated with
a simulation-variable before the initialization of the fault.

\subsubsection {Random Number Generation}\label{sec:UG:random_node}
Fault Management has its own class for generating random numbers, called
\verb"FaultRandNumber".
This class is used in some types of faults to generate noise and in the fault
manager to set random trigger values and fault parameters.

\par The user can configure an \verb"FaultRandNumber" directly or through the
attributes of a node in XML called \verb"RandValue":
\begin{center}\begin{tabular}{|p{3.1cm}|p{2.4cm}|p{8.5cm}|} \hline
\textbf{RandNumber Member} & \textbf{RandValue Attribute} & \textbf{Description}
\\ \hline
distribution\_type & distribution & The type of distribution. Can be
\verb"GAUSSIAN" (Normal) or \verb"FLAT" (Uniform). \\ \hline
mean & mean & The mean of the Gaussian distribution. Defaults to 0. \\ \hline
std\_dev & std\_dev & The standard deviation of the Gaussian distribution. \\
\hline
lower\_limit & min & The lower limit of the uniform distribution. \\ \hline
upper\_limit & max & The upper limit of the uniform distribution. \\ \hline
seed & seed & The seed for the random number generator. \textbf{
              If this is not provided, OR SET TO 0, the model uses the
              system time to generate a seed}. \\ \hline
\end{tabular}\end{center}

This random number generation capability is used in the following places:
\begin{itemize}
\item For setting the (per-cycle) value associated with the white-noise and
      random-walk Faults. The white-noise and random-walk Faults have an
      instance of \verb'FaultRandNumber' built into their classes;
      \textit{FaultWhiteNoise::noise}, and \textit{FaultRandomWalk::rand}.
      These two Faults will
      generate a new random value according to their assigned random
      distributions on every cycle for which the Fault is triggered.
\item For setting the (static) value of a random overwrite Fault;
      this value will be applied consistently on every cycle for which the
      Fault is triggered, but with a different seed will generate a different
      value on a different run. Because this type uses a single random number
      per run, it does not need its own \verb'FaultRandNumber' instance and
      can use one from elsewhere.
\item For generating the value at which a Trigger triggers.
      Like the overwrite-fault, because this use-case uses a single random
      number per run, it does not need its own \verb'FaultRandNumber'
      instance and can use one from elsewhere.
\end{itemize}





\subsection{Configuration Settings}
The various components are configured differently depending on which
architecture is used to implement the fault-management model. Details follow in
these specific sections.





\section{Architecture: Direct Instantiation} \label{sec:UG:DirectInstantiation}
This architecture is well-suited to relatively simple fault-injection
capabilities, and for quick testing. For more complex systems, the fault-manager
would be better suited (see Section \ref{sec:UG:FaultManager}).

\subsection{Top-level Instantiation and Methods}
Faults can be instantiated directly into the code representing the operation of
some sensor, and configured to introduce their respective errors inline with the
model code. In the following simple example,
we want to apply a bias-fault to a working variable in the processing of
the sensor, and a scale-fault to the sensor output at the end of the
processing.

We need a trigger and trigger-group, and instances of the specific faults being
applied:
\begin{codeC}
#include "cml/models/utilities/fault_management/include/fault_bias.hh"
#include "cml/models/utilities/fault_management/include/fault_scale.hh"
#include "cml/models/utilities/fault_management/include/trigger.hh"
#include "cml/models/utilities/fault_management/include/trigger_group.hh"

class Sensor {
 public:
  double working_variable; /* (--) some internal value in the class. */
  double output; /* (--) the model output. */
  Trigger<double> trigger;
  TriggerGroup trigger_group;
  FaultBias<double> early_fault;
  FaultScale<double> later_fault;
\end{codeC}

The trigger and the fault must be passed their variables during construction.
Note that the template argument of these classes must match the type of the
variable being passed in.
\begin{codeC}
Sensor(double& sim_time) :
  trigger(sim_time), // sim_time is the trigger variable.
  early_fault(working_variable),
  later_fault(output)
{}
\end{codeC}

Each trigger group (in this example there is only 1) must be populated with its
associated triggers; each trigger-group can have multiple triggers, and each
trigger can be associated with multiple trigger-groups.

Each fault must also be populated with its associated trigger-groups; each fault
can have multiple groups, and each group can be associated with multiple faults.

The corresponding methods -- \textit{add\_trigger} and
\textit{add\_trigger\_group} -- may be called from the input file or from part
of the sensor initialization routine.

Each fault must be initialized \textbf{after} adding the trigger groups.
\begin{codeC}
void Sensor::initialize() {
  trigger_group.add_trigger(&trigger);
  early_fault.add_trigger_group(&trigger_group);
  later_fault.add_trigger_group(&trigger_group);
  early_fault.initialize();
  later_fault.initialize();
}
\end{codeC}

In the sensor-model update-code, we can now insert the points at which we want
to apply faults.
\begin{codeC}
void Sensor::update() {
  // Preliminary sensor processing
  // working_variable = some function of inputs
  ...
  // Inject the fault (if the trigger condition is met) on working_variable
    early_fault.update();

  // More sensor processing
  // output is generated as a function of the biased working_variable
  ...
  // Inject the second fault to output (if the trigger condition is met).
  later_fault.update();
}
\end{codeC}

Note -- in this example, we used \verb"FaultBias" and \verb"FaultScale" classes
for our selected faults. Other classes are available for different types of
faults.
See Section \ref{sec:UG:components:faults} for a table listing the fault types
and classes and for details on how to configure them. Examples of
instantions follow (e.g. these might add to the instances found in the
definition of the Sensor class above, then use similar construction
requriements).
\begin{codeC}
  FaultBias<double> bias_fault;
  FaultScale<double> scale_fault;
  FaultOverwrite<double> overwrite_fault;
  FaultStale<double> stale_fault;
  FaultWhiteNoise<double> white_noise_fault;
  FaultRandomWalk<double> random_walk_fault;
  FaultFunction<double> function_fault;
\end{codeC}

\subsection{Independent Variables}
Independent variables are used in function-faults; these must be associated with
a simulation-variable before initializing the fault. Because of limitations of
passing references via SWIG, it is difficult to make this connection via the
input file; it is recommended that this association be defined in code, usually
in the \textit{Sensor::initialize()} method.

For example, suppose \textit{function\_fault} was included as an instance
in the \textit{Sensor} class defintion, then somewhere in
\textit{Sensor::initialize()} we might expect to see:
\begin{codeC}
function_fault.ind_variable.initialize(independent_variable);
function_fault.initialize();
\end{codeC}
Note that this should be handled in compiled code rather than the Python
input-processor because it relies on templates.

Independent variables default to a \textit{relative} behavior by default
(i.e. the value used is the value relative to the reference value recorded at
triggering); they can be made \textit{absolute} by
setting the \textit{relative\_value} flag to false.
See section \ref{sec:formulation:indep_var} for discussion of absolute
versus relative.

This can be assigned in
either the same coded method, or through the input file.
\begin{codePython}
 vehicle.sensor.function_fault.relative_value = False
\end{codePython}


\subsection{ Configuration -- Triggers}
\label{sec:UG:DI:config_trigger}
The trigger-variable must be assigned to the trigger during the
\textbf{construction} of
each trigger; the operator and value should be set either during
initialization or in an input file.
\begin{codePython}
vehicle.sensor.trigger_A.Operator = trick.TriggerBase.GE
vehicle.sensor.trigger_A.value = 50.0
\end{codePython}
It is also recommended to give each trigger a unique name.
This allows triggers to be disabled within specific trigger groups.
\begin{codePython}
vehicle.sensor.trigger_A.name = "trigger A"
\end{codePython}

The enabled status of a trigger is specific to each trigger-group with which it
is associated -- it can be enabled in one and disabled in another. This setting
is modified with the \textit{set\_trigger\_enable} method:
\begin{codePython}
vehicle.sensor.trigger_group_X.set_trigger_enable("trigger A", False)
\end{codePython}

To make a trigger periodic (see Section \ref{intro_periodic}), use the
\textit{set\_periodic} method.
\begin{codePython}
 vehicle.sensor.trigger_A.set_periodic(2.0, 5.0)
\end{codePython}

This can be undone by calling \textit{unset\_periodic}.
\begin{codePython}
 vehicle.sensor.trigger_A.unset_periodic()
\end{codePython}

To limit the number of consecutive times a Trigger can be used, use the
\textit{set\_trigger\_count( number)} function.
\begin{codePython}
 vehicle.sensor.trigger_A.set_trigger_count( 5)
\end{codePython}

This can be undone by calling \textit{unset\_trigger\_count}.
\begin{codePython}
 vehicle.sensor.trigger_A.unset_trigger_count()
\end{codePython}

\subsection{Configuration - Faults}
The variable to be modified by the fault is passed in during construction.

\subsubsection{Parameters}
\label{sec:UG:DI:config_fault:param}
Each type of fault (except stale-value) has additional parameters specific to
that type.
Most of them can be modified in input files through the fault's
\textit{set\_param} method, whose arguments are the name of the parameter (as a
string) and the value to set the parameter to. Note that several sub-classes of
the base class \textit{FaultFunctionBase} add new parameters; the
\textit{set\_param} method for these sub-classes is extended accordingly.
\begin{codePython}
 vehicle.sensor.white_noise_fault.set_param(`max',10.0)
\end{codePython}
Alternatively, \textbf{some of} these parameters may be set by direct
assignment, if the corresponding variable name is known:
\begin{codePython}
 vehicle.sensor.white_noise_fault.noise.max = 10.0
\end{codePython}
The difference between these techniques is one of safety. Different
fault-types use different parameters, so if attempting tos et variables
directly, it is important to confirm that the target fault does have the
corresponding parameter. For example, a fault of type \textit{FaultOverwrite}
has a \textit{"value"} parameter setting, internally identified as
\textit{faulted\_value}. If an invalid
parameter is assigned while using \textit{set\_param}, an error message will be
posted, the assignment ignored, and the simulation will proceed.
If an invalid parameter is assigned
directly from the input file, the simulation will terminate.

These two assignments are valid and equivalent:
\begin{codePython}
 vehicle.sensor.white_noise_fault.set_param(`max',10.0)
 vehicle.sensor.white_noise_fault.noise.max = 10.0
\end{codePython}
as are these two:
\begin{codePython}
 vehicle.sensor.overwrite_fault.set_param(`value',10.0)
 vehicle.sensor.overwrite_fault.faulted_value = 10.0
\end{codePython}

These assignments are invalid because these particular faults have no
parameters by these names. Because these assignments are made using
\textit{set\_param(...)}, they will result in an error message and be
otherwise ignored.
\begin{codePython}
 vehicle.sensor.white_noise_fault.set_param(`value',10.0)
 vehicle.sensor.overwrite_fault.set_param(`max',10.0)
\end{codePython}

These direct assignments are invalid because they use the wrong variable
name for the target parameter. Because they are made by direct assignments,
these will result in simulation termination:
\begin{codePython}
 vehicle.sensor.white_noise_fault.max = 10.0
 vehicle.sensor.overwrite_fault.value = 10.0
\end{codePython}
These direct assignments are invalid because they use a parameter name that
is valid only for a different type of fault, and is invalid for the
specified fault type. These, too, will result in simulation termination:
\begin{codePython}
 vehicle.sensor.white_noise_fault.faulted_value = 10.0
 vehicle.sensor.overwrite_fault.noise.max       = 10.0
\end{codePython}



\subsubsection{Example - Linear Function Faults}
Linear functions are defined by a rate and a nominal value.
These represent the values $m$ and $b$ respectively in $y=mx+b$.
\begin{codePython}
vehicle.sensor.function_fault.type = trick.FaultFunctionBase.Linear
vehicle.sensor.function_fault.nominal = 2.0
vehicle.sensor.function_fault.rate = -0.05
\end{codePython}
Or, equivalently:
\begin{codePython}
vehicle.sensor.function_fault.type = trick.FaultFunctionBase.Linear
vehicle.sensor.function_fault.set_param('nominal',2.0)
vehicle.sensor.function_fault.set_param('rate', -0.05)
\end{codePython}

\subsubsection{Periodic Function Faults}
Three periodic functions are implemented: sine waves, square waves, and
triangle waves.
These are defined by a frequency, an amplitude, and a phase-offset.
Each of these three parameters can be represented as a linear equation of some
independent variable, using parameters \verb"nominal" and \verb'rate'.  If only
the \verb"nominal" parameter is set for any of these three parameter, the
respective parameter will be constant at that value.
Alternatively, by defining a \verb"rate" and initializing the independent
variable, these three controlling parameters become linear functions themselves.
In the example below, we use a sine-wave function with constant frequency, an
amplitude that starts at value 5.0 and ramps down linearly with \verb"sim_time",
and a phase-offset that starts at 0.0 and ramps up linearly with another
variable, \verb"other_variable".
\begin{codeC}
// Somewhere in Sensor::initialize()
function_fault.ind_variable.initialize(sim_time);
function_fault.amplitude.ind_variable.initialize(sim_time);
function_fault.phase_offset.ind_variable.initialize(other_variable);
function_fault.initialize();
\end{codeC}
Note that the initialization should be handled in compiled code rather than the
Python input-processor because it relies on templates.

The parameter values may be set in code or in the input-processor
\begin{codePython}
sensor_object.function_fault.type = trick.FaultFunctionBase.Sinewave
sensor_object.function_fault.frequency.nominal = 2.0
sensor_object.function_fault.amplitude.nominal = 5.0
sensor_object.function_fault.amplitude.rate = -0.1
sensor_object.function_fault.phase_offset.rate = 1.0
\end{codePython}

\subsubsection{Fire Limit}
To set a fire limit (see Section
\ref{sec:formulation:fault_support:fire_count}), set the
\textit{is\_fire\_limited}
flag to true and, if desired, set the number of times the fault will fire,
\textit{fire\_limit} (the default is 1 frame).
\begin{codePython}
 vehicle.sensor.scale_fault.is_fire_limited = True
 vehicle.sensor.scale_fault.fire_limit = 10
\end{codePython}

\subsubsection{Enabled} \label{sec:UG:DirectInstantiation:Config:Enabled}
To enable/disable the fault, set the \verb"enabled" flag accordingly.
\textbf{Faults are disabled by default.}
\begin{codePython}
 vehicle.sensor.scale_fault.enabled = True
\end{codePython}

The \verb'enabled' flag is public, so it is possible to disable and enable a
fault from the input file. There is also a \textit{disable()} method that may be
used to disable a fault; there is a subtle difference between setting
\textit{enabled = False} and calling \textit{disable()}:
\begin{itemize}
 \item Disabling and later re-enabling the fault by directly setting
   \textit{enabled} \textbf{will not} automatically reset the fault, which might
   be the expected behaviour if the fault is in the triggered state when it is
   both disabled and re-enabled.
  \item calling \textit{disable()} and later (or immediately) re-enabling the
    fault \textbf{will} result in the fault resetting.
  \item The difference between these two operations lies in the logic that
    triggers a reset -- a fault gets reset if it is in a triggered state with
    the protected flag \textit{was\_triggered\_last\_pass} set to false. Simply
    setting enabled to false does not (un)set
    \textit{was\_triggered\_last\_pass},
    but the \textit{disable()} method does assign
    \textit{was\_triggered\_last\_pass = false}.
\end{itemize}

Note that if the fault is fire-limited, it will disable itself upon reaching the
fire-count with the \textit{disable()} method, so will reset after being
re-enabled.

\subsection {Random Number Generation}
The random-walk and white-noise faults have built-in instances of the model's
random-number generation capability (see Section \ref{sec:UG:random_node}.
These can be accessed directly to support
the \verb'FaultRandNumber' methods \textit{initialize\_gaussian} and
\textit{initialize\_flat}. The \verb'FaultRandNumber' instances are found at
\textit{FaultWhiteNoise::noise} and \textit{FaultRandomWalk::rand}. Both
methods take 2 or 3 arguments. In both cases the 3rd argument is optional
and specifies the seed for the random number generator; if not specified,
(or if specified to 0) the model will generate a seed.
If reproducible results are desired, the 3rd
argument is required (and be non-zero) because the generated seed will not
be the same next time:
\begin{itemize}
\item \textit{initialize\_gaussian} {mean, standard-deviation, seed}
\item \textit{initialize\_flat} {minimum, maximum, seed}
\end{itemize}
For example:
\begin{codePython}
vehicle.sensor.white_noise_fault.noise.initialize_gaussian(0.0, 2.0)
vehicle.sensor.random_walk_fault.rand.initialize_flat(0.0, 2.0,5)
\end{codePython}

To randomize a trigger value or an overwrite-fault value (as discussed in
Section \ref{sec:UG:random_node}), any random engine
can be used to generate a random number and directly assign it to the
respective
Trigger's \textit{value} (see Section \ref{sec:UG:DI:config_trigger}) or the
overwrite-fault's \textit{faulted\_value} (see Section
\ref{sec:UG:DI:config_fault:param}) variable.


\section{Architecture: Fault Manager} \label{sec:UG:FaultManager}
For more complicated situations, it may be desirable to use the fault-manager
design and the configure the faults with an XML file. This has the advantage of
being a single instantiation, and the disadvantage of relying on dynamic memory
allocation via the \textit{new/delete} process.
This design pattern still requires that the sensor code have faults inlined into
the model's standard execution, but the syntax for configuring these faults
and for inserting them into the code
is substantially different because the faults cannot be addressed directly.

\subsection{Top-level Instantiation and Methods}
\par The instantiation of a fault manager is handled within the class
representing the sensor being faulted, but now we do not need triggers, trigger-
groups, or specific faults.
\begin{codeC}
#include "cml/models/utilities/fault_management/include/FaultManager.hh"

class Sensor {
 public:
  ...
  FaultManager faults;
\end{codeC}

In the direct-instantiation architecture,
the sensor model had direct access to each fault and could call them as needed,
but that flexibility is not so easily available when using
the fault-manager architecture. Instead, the fault-manager tags each of its
faults with one of five ``locations'', breaking the faults into five sets, with
each set applied at a specific point in the update cycle of the parent / sensor
model. In this regard, the direct-instantiation architecture provides more
flexibility because faults may be injected at any number of locations within
the model, while the fault-manager architecture limits the locations to 5.
In reality, this is not a significant issuel there are few algorithms that
have more than 5 meaningful locations at which a fault could be applied.
The fault-manager \textit{update()}
method requires a single argument to identify which set of faults is to be
applied. The five sets are:
\begin{itemize}
 \item \textit{\textbf{Init:}} \verb'INIT': faults in this set are applied at
   initialization.
 \item \textit{\textbf{Upstream:}} \verb'UPSTREAM': faults in this set are
   applied at the top of a frame.
 \item \textit{\textbf{Intermediate\_1:}} \verb'INTERMEDIATE_1': faults in this
   set are applied part-way through the update to the sensor state.
 \item \textit{\textbf{Intermediate\_2:}} \verb'INTERMEDIATE_2': faults in this
   set are applied part-way through the update to the sensor state; these are
   applied after the \textit{Intermediate\_1} set.
 \item \textit{\textbf{Downstream:}} \verb'DOWNSTREAM': faults in this set are
   typically applied after the sensor model has completed generating its output.
 \item (Additional) \textit{\textbf{INVALID:}} is an error state, used for
   debugging.
\end{itemize}


The fault-manager is initialized along with the sensor model whose values it is
modifying. Any initialization-specific faults can be applied by calling the
fault-manager \textit{update()} method with the \verb'Init' location.
\begin{codeC}
void initialize() {
  faults.initialize();
  ...
  faults.update(FaultManager::Location::Init);
}
\end{codeC}

The fault-manager's \textit{update()} method is called from the appropriate
locations in the the sensor model's \textit{update()} method with the
appropriate tag to identify which set of faults to apply.
\begin{codeC}
  void update() {
    faults.update(FaultManager::Location::Upstream);
    // Sensor processing
    ...
    faults.update(FaultManager::Location::Intermediate_1);
    // More sensor processing
    ...
    faults.update(FaultManager::Location::Intermediate_2);
    // More sensor processing
    ...
    faults.update(FaultManager::Location::Downstream);
  }
}
\end{codeC}

The fault manager is disabled by default; it must be enabled by setting a flag
and providing the name of an XML file describing the faults.
This is generally done in an input file.
\begin{codePython}
vehicle.sensor.faults.enabled = True
vehicle.sensor.faults.fault_file = "fault_file.xml"
\end{codePython}

Note -- \textit{vehicle.sensor} is an instance of the Sensor class just defined.

\subsection{Configuration - XML file}

\subsubsection{XML Structure Primer}
An XML file defines a set of nodes.
Each node has a name, any number of attributes, and any number of child nodes.
There is no limit to how many nodes or levels of nodes are defined in the file,
but the top level must contain a single root node.

Throughout this document, we will refer to \textbf{node-names} and
\textbf{field-names}.
\begin{itemize}
\item node-name refers to the type of node; in the example below, this might
      be \textit{Node1}, \textit{ChildNode1},
      \textit{GrandchildNode1}, etc.
\item field-name refers to the identifier associated with some attribute.
      in the example below, this might be
      \textit{child\_attribute} as a field-name of the \textit{ChildNode1}
      node, or \textit{grandchild\_att} as a fileld-name of the
      \textit{GrandchildNode1} node; the \textit{GrandchildNode2} node has
      no field-names.
\end{itemize}


The general XML format looks like:
\begin{codeXML}
<?xml version="1.0" encoding="UTF-8">
<RootNodeName root_attribute="">
<Node1 attribute1="value1" attribute2="value2">
  <ChildNode1 child_attribute="child value">
    <GrandchildNode1 grandchild_att="another value"></GrandchildNode1>
    <GrandchildNode2></GrandchildNode2>
  </ChildNode1>
  <ChildNode2></ChildNode2>
</Node1>
<Node2>
  <!-- Et cetera. By the way, this is what a comment looks like. -->
</Node2>
</RootNodeName>
\end{codeXML}

\subsubsection{Fault Management XML Structure}
A fault management XML file must have:
\begin{itemize}
\item a root node called \verb"FaultManager".
\begin{codeXML}
<?xml version="1.0" encoding="UTF-8"?>
<FaultManager name="Test Faults">
\end{codeXML}
\item which has one or more \verb'Fault' sub-nodes
  \begin{codeXML}
  <Fault ...>
  ...
  </Fault>
\end{codeXML}
\item Each \verb'Fault' node contains:
  \begin{itemize}
    \item A \verb'SimVar' sub-node specifying the simulation variable associated
     with the fault:
    \begin{codeXML}
    <SimVar name="test_object.var_char"></SimVar>
    \end{codeXML}
    \item One or more \verb'TriggerGroup' sub-nodes identifying the trigger-
      groups associated with the fault.
    \begin{codeXML}
    <TriggerGroup>
    ...
    </TriggerGroup>
    \end{codeXML}
    Each \verb'TriggerGroup' node contains one or more Trigger sub-nodes
    identifying the triggers associated with the trigger-group
      \begin{codeXML}
      <Trigger ...> </Trigger>
      \end{codeXML}
  \end{itemize}
\end{itemize}

We will look at the structure of each of these node types, starting with
Section \ref{sec:UG:FM:XML_config_fault_node}.

\subsubsection{Case Sensitivity}
\label{sec:UG:FM:Config_xml:case}
In general the node-names and field-names as specified in this User's Guide
are case-sensitive. The model will not find a node or field if its appearance
in the XML file does not match with its expected name specified in this
document. For example, we discuss the \textit{Fault} node in several places.
The model is expecting to find XML content of the form
\begin{codeXML}
<Fault ...></Fault>
\end{codeXML}

The model will not recognize:
\begin{codeXML}
<fault ...></fault>
\end{codeXML}
or
\begin{codeXML}
<FAULT ...></FAULT>
\end{codeXML}

The only exception to this is for the field-name \textit{type}, found in a
\textit{Fault} node or \textit{Function} node. For both cases, the
field-name \textit{Type} will also be recognized. The XML-helper model, which is
used by this model to parse the XML files, does have the option to ignore
case on the first character of a field-name. This Fault Management model
exercises that
option for this  \textit{type / Type} field-name only; this is done to
support some legacy inconsistencies. At some point, if it would be
beneficial to do so, we could allow other
field-name lookups to utilize the same option.

Note -- this option requires editing the \textit{fault\_manager.cc} file
to add a \textit{true} flag as the (optional) third argument in the
respective \textit{XmlHelper::xml\_find\_value(...)} calls.

\subsection {XML Configuration -- Random Values}
As discussed in Section \ref{sec:UG:random_node}, random values may be used
in different places throughout the configuration. Random distributions, from
which the random values are drawn, are typically configured in a
\textit{RandValue} node. A \textit{RandValue} node requires:
\begin{itemize}
\item a \textit{distribution} field, specifying what type of random
      distribution to use. The model currently supports Normal distributions
      (\textit{GAUSSIAN}), and Uniform distributions (\textit{FLAT}).
\item values describing the characteristic shape of the distribution.
\item (optional) a \textit{seed} value for seeding the random-number
      generator. If this field is omitted, \textbf{OR if it is set to 0}, the
      model will generate a seed from the system clock. While this field is
      technically optional, if the results are to be repeatable, it is
      required.`
\end{itemize}

\subsubsection {Normal Distribution}
A Normal distribution is configured with the following fields:
\begin{itemize}
\item \textit{distribution="GAUSSIAN"}
\item \textit{std\_dev} to specify the standard deviation.
\item \textit{mean} to specify the mean value. This field is optional, it
      defaults to 0.0 if not specified.
\end{itemize}
The model uses the \textit{normal\_distribution} from the C++
\textit{<random>} library. This implementation expectes that the standard
deviation be $>0$. The model may run if this is not
satisfied, but it is not guaranteed.

\subsubsection {Uniform Distribution}
A Uniform distribution is configured with the following fields:
\begin{itemize}
\item \textit{distribution="GAUSSIAN"}
\item minimum/maximum specifications using either:
  \begin{itemize}
  \item \textit{min}
  \item \textit{max}
  \end{itemize}
  or
  \begin{itemize}
  \item \textit{rel\_min}
  \item \textit{mean}
  \item \textit{rel\_max}
  \end{itemize}
\end{itemize}
With the first option (min/max), the distribution generates a value
$x \in [min, max)$.
Both fields must be specified to use this option. This is the default
option; if both fields are specified, the model will not parse the contents
for the second option.

With the second option, the distribution generates a value
$x \in [mean - |rel\_min|, mean + rel\_max)$. Note that \textit{rel\_min}
can be specified as positive or negative; its absolute-value will be
subtracted from the \textit{mean}. All three fields must be specified to use
this option.

The model uses the \textit{real\_uniform\_distribution} from the C++
\textit{<random>} library. This implementation expectes that the lower bound
will be less than the upper bound. The model may run if this is not
satisfied, but it is not guaranteed.

\begin{codeXML}
<RandValue><distribution="GAUSSIAN" mean="1.0" std_dev="2.0"></RandValue>
...
<RandValue><distribution="FLAT" min="-1.0" max="1.0" seed="1"></RandValue>
\end{codeXML}


\subsection {XML Configuration -- Fault Nodes (Common)}
\label{sec:UG:FM:XML_config_fault_node}
Fault nodes contain the following attributes:
\begin{itemize}
 \item \verb'ID' is a name for the fault
 \item \verb'enabled' sets the enabled status of this fault. If not specified,
   this defaults to false.
 \item \verb'type' identifies the type of fault -- e.g. STALE, FUNCTION.
   Note -- \verb'Type' is equivalent to \verb'type' (see Section
   \ref{sec:UG:FM:Config_xml:case}).
 \item \verb'loc' identifies the location of where the fault will be applied.
\end{itemize}
\begin{codeXML}
<Fault ID="Bias Fault" type="BIAS" loc="DOWNSTREAM">
  ...
</Fault>
\end{codeXML}

For each type of fault-node, there may be specific additional configurations;
these are explored below.

\subsection{XML Configuration -- Bias Fault Nodes}
\par There is one parameter specific to bias faults: the bias.
This is the \verb"bias" member of the class.
It is specified by field-name \verb'value' in a \verb'Bias' node.
\begin{codeXML}
<Fault ID="Bias Fault" type="BIAS" loc="DOWNSTREAM">
  <Bias value="100"></Bias>
  ...
</Fault>
\end{codeXML}

\subsection{XML Configuration -- Scale Fault Nodes}
\par There is one parameter specific to scale faults: the scale factor.
This is the \verb"scale_factor" member of the class.
It is specified by field-name \verb'value' in a \verb'Scale' node.
\begin{codeXML}
<Fault ID="Scale Fault" type="SCALE" loc="DOWNSTREAM">
  <Scale value="10"></Scale>
  ...
</Fault>
\end{codeXML}

\subsection{XML Configuration -- Overwrite Fault Nodes}
\par There is one parameter specific to overwrite faults: the overwrite value.
This is the \linebreak \verb"overwrite_value" member of the class.
It is usually specified by field-name \verb'value' in a \verb'Overwrite' node.

\begin{codeXML}
<Fault ID="Overwrite Fault" type="OVERWRITE" loc="DOWNSTREAM">
  <Overwrite value="100"></Overwrite>
  ...
</Fault>
\end{codeXML}

The overwrite value may be also be configured to be a randomized value;
the distribution for the random number generation is
set up using a RandValue node (see Section \ref{sec:UG:random_node}).
\begin{codeXML}
<Fault ID="Random Overwrite Fault" type="OVERWRITE" loc="DOWNSTREAM">
  <RandValue distribution="GAUSSIAN" mean="150" std_dev="10" seed="1">
  </RandValue>
</Fault>
\end{codeXML}

\subsection{XML Configuration -- White-Noise and Random-Walk Fault Nodes}
\par The noise in white-noise and random-walk faults is generated by
\verb"FaultRandNumber" instances, configured in \verb'RandValue' nodes.
See Section \ref{sec:UG:random_node} for information on how to configure them.
\begin{codeXML}
<Fault ID="White-Noise Fault" type="WHITENOISE" loc="DOWNSTREAM">
  <RandValue distribution="GAUSSIAN" std_dev="5" seed="1"></RandValue>
  ...
</Fault>
\end{codeXML}

\subsection{XML Configuration -- Function Fault Nodes}
When a function fault is triggered, it biases its variable by some function
of some independent variable. This type of fault is identified by the
\verb'type="FUNCTION"', and the specific type of function is specified by the
\verb"type" member of the \verb'Function' sub-node, using an enumeration from
\verb"FunctionType". Available function types are:
\begin{itemize}
 \item LINEAR
 \item SINEWAVE
 \item TRIANGLEWAVE
 \item SQUAREWAVE
\end{itemize}

\begin{codeXML}
  <Fault ID=... type="FUNCTION" loc= ...>
    ...
    <Function type="LINEAR">
    ...
    </Function>
  </Fault>
\end{codeXML}

\subsubsection{Independent Variables}
Independent variables must be initialized with a variable reference. In XML
files, independent variables are represented by a \verb"IndVariable" node with
\verb"name" and \verb"relative" attributes, embedded as a sub-node in the
\verb'Function' node.
The \verb"relative" attribute is optional (default is true).
Like a trigger variable or fault variable, an independent variable can be any
variable in the simulation, even if it is private or const, as long as the type
is valid.
\begin{codeXML}
  <Fault ID=... type="FUNCTION" loc= ...>
    ...
    <Function type=...>
      <IndVariable name="time.sim_seconds" relative="true"></IndVariable>
      ...
    </Function>
    ...
  </Fault>
\end{codeXML}

\subsubsection{Parameters}
The \textit{nominal} and \textit{rate} variables within a function-fault are
set with a \verb'Parameters' sub-node of the \verb'Function' node. Note that
these parameters are only used for a linear function fault.

The two values are the $m$ (\textit{rate}) and $b$ (\textit{nominal}) in
$y=mx+b$.


\begin{codeXML}
 <Fault ID=... type="FUNCTION" loc= ...>
    ...
    <Function type="LINEAR"...>
      <IndVariable ...></IndVariable>
      <Parameters nominal="2.0" rate="-0.05"></Parameters>
    </Function>
    ...
</Fault>
\end{codeXML}

The \textit{initial} keyword can also be used as a substitute for
\textit{nominal} to represent the y-intercept. Neither is a particularly good
name -- nominal might be interpreted as "no-fault", and initial might be
interpreted as the first application of the fault, which is only valid if the
first application has $x=0$.

\begin{codeXML}
      ...
      <Parameters initial="2.0" rate="-0.05"></Parameters>
      ...
\end{codeXML}
If \textit{nominal} and \textit{initial} are both specified, \textit{nominal}
takes precedence. If neither is specitied, it defaults to 0.0 and the linear
equation becomes $y=mx$. Note that specification of \textit{rate} is required.

\subsection{XML Configuration -- Periodic Function Fault Nodes}
Three periodic functions are implemented: sine waves, square waves, and
triangle waves.
These are defined by a frequency, an amplitude, and a phase-offset, represented
as separate sub-nodes of the \verb'Function' node.
Each of these three parameters can be represented as a linear equation of some
independent variable, using parameters \verb"nominal" and \verb'rate', and sub-
node \verb'IndVariable'.
\begin{itemize}
 \item If only the \verb"nominal" parameter is set for any of these three
 parameters, the respective parameter will be constant at that value.
 \item If the \verb"nominal" parameter is not set, it will default to 0.0.
 \item If one of these nodes contains a \verb"rate" attribute but not an
 \verb"IndVariable" sub-node, it will use the same independent variable as the
 function itself.
\end{itemize}

In the example below, we use a sine-wave function with constant frequency; an
amplitude that starts at value 5.0 and ramps down linearly with \verb"sim_time",
the same variable used as the independent variable for the sine-wave itself;
and a phase-offset that starts at 0.0 and ramps up linearly with another
variable, \verb"other_variable".

\begin{codeXML}
<Fault ...>
  ...
  <Function type="SINEWAVE">
    <IndVariable name="sim_time" relative="true"></IndVariable>
    <Frequency nominal="2.0"></Frequency>
    <Amplitude nominal="5.0" rate="-0.1"></Amplitude>
    <PhaseOffset rate="1.0">
      <IndVariable name="sim_object.other_variable"></IndVariable>
    </PhaseOffset>
  </Function>
  ...
</Fault>
\end{codeXML}

\subsection {XML Configuration -- SimVar Nodes}
The SimVar node has a single attribute, \verb'name', which specifies the
simulation address of the faulted variable.
\begin{codeXML}
<Fault ...>
  <SimVar name="test_object.var_char"></SimVar>
  ...
</Fault>
\end{codeXML}

\subsection {XML Configuration -- TriggerGroup Nodes}
Trigger-groups are unique to a fault. They may be replicated, but not re-used.
Trigger-group nodes comprise only a set of Trigger sub-nodes and no additional
attitubutes.
\begin{codeXML}
<Fault ...>
  ...
  <TriggerGroup>
    ...
  </TriggerGroup>
</Fault>
\end{codeXML}

\subsection {XML Configuration -- Trigger Nodes}
Trigger nodes are embedded within the trigger-group node for a specific fault.
The first time a trigger appears in the file, its node must have the following
attributes:
\begin{itemize}
 \item \verb"name" should be unique; this name can be used later to associate
   this trigger with another trigger-group.
 \item \verb"variable" is the simulation address of the trigger-variable (i.e.
   the variable whose value is being monitored for triggering the fault). Note
   that the variable can be any variable in the simulation, even if it is
   private or const, as long as the data-type is supported.
 \item \verb"compare" is the operator being used to compare the value of the
   trigger variable against the trigger value threshold.
 \item Either:
 \begin{itemize}
  \item \verb"value" provides the trigger-value, or
  \item a \verb"RandValue" sub-node (see Section \ref{sec:UG:random_node}), used
    if the trigger-value is determined by a random variable.
 \end{itemize}
 \item (optional) A Trigger can be made periodic (see Section
   \ref{intro_periodic}) by adding the
   \verb'trigger_period' and \verb'trigger_length' attributes.
 \item (optional) A Trigger can be limited to triggering no more than some
   number of consecutive times (see Section
   \ref{sec:formulation:fault_support:trigger_fire_count}
   by adding the
   \verb'fire_limit' attribute.
\end{itemize}
\begin{codeXML}
<Fault ...>
  ...
  <TriggerGroup>
    <Trigger name="trigger XYZ" variable="sim_object.sim_variable"
      compare="GE" value="5"></Trigger>
    ...
  </TriggerGroup>
</Fault>
\end{codeXML}

When a trigger is being used in another trigger-group, the syntax is simpler,
containing only one attribute:
\begin{itemize}
 \item \verb"reuse_name" identifies which of the previously defined
   trigger-nodes to use in this trigger-group.
\end{itemize}
\begin{codeXML}
<Fault ...>
  ...
  <TriggerGroup>
    <Trigger reuse_name="trigger_XYZ"></Trigger>
    ...
  </TriggerGroup>
</Fault>
\end{codeXML}


\subsection{Example}
\begin{codeXML}
<Fault ID="Discrete Noise" type="FUNCTION" loc="DOWNSTREAM">>
  <SimVar name="vehicle.sensor_X.output"></SimVar>
  <Function type="SQUAREWAVE">
    <IndVariable name="time.sim_time" relative="true"></IndVariable>
    <Frequency nominal="2.0"></Frequency>
    <Amplitude nominal="5.0" rate="-0.1"></Amplitude>
    <PhaseOffset rate="1.0">
      <IndVariable name="sim_object.other_variable"></IndVariable>
    </PhaseOffset>
  </Function>
  <TriggerGroup>
    <Trigger name="trigger ABC" variable="sim_object.sim_variable"
      compare="GE" value="5" trigger_period="2" trigger_length="0.5"
      fire_limit="10">
    </Trigger>
    <Trigger reuse_name="trigger_XYZ"></Trigger>
  </TriggerGroup>
</Fault>
\end{codeXML}




\subsection{Post-XML Modifications}
With the direct-instantiation architecture, modifying a component of the
fault-model after initialization is straightforward because we have direct
access to
each of the components. With the fault-management architecture, we need to gain
access via the component's name and the fault-manager's lookup methods:
\begin{itemize}
 \item \verb'Fault * get_fault( name)'
   returns a pointer to the fault with the specified name, or ID.
 \item \verb'TriggerBase * get_trigger(name)'
   returns a pointer to the trigger with the specified name.
 \item \verb'set_fault_enabled(fault_name, enable_flag)'
   sets the enabled flag on the fault with the specified name.
 \item \verb'set_fault_trigger_enabled(fault_name, trigger_name, enable_flag)'
   sets the
   enabled status of the specified trigger in all trigger groups the specified
   fault. Recall, a trigger can be enabled independently per trigger-group for
   each trigger-group in which it is embedded; this method sets the enabled flag
   for the specified trigger for all trigger-groups in the fault.
 \item \verb'set_fault_param( fault_name, param_name, double value)'
   sets the parameter with the specified name in the fault with the speciied
   name to the specified
   value.
 \item \verb'set_trigger_value( trigger_name, double value)' sets the trigger-
   value for the specified trigger.
\end{itemize}

\textbf{Important Note} -- The use of \textit{set\_fault\_param} and
\textit{set\_trigger\_value} in the input file, to be processed with the
rest of the input file, will take precedence over values specified in the
XML file. It is \textbf{NOT} necessary to delay modifications to
configurations introduced via the input file until after
the simulation has been fully initialized. When these methods are executed
before the XML file is parsed, the settings are cached and applied at
initialization (after parsing the XML file and applying its settings).








\chapter{Verification}
\par Fault Management is verified by two unit simulations: \textit{SIM\_fault}
and \textit{SIM\_fault\_manager}.
In \textit{SIM\_fault\_manager}, the faults are all implemented and
configured via the XML file interface following the patterns described in
Section \ref{sec:UG:FaultManager}.
In \textit{SIM\_fault} the faults are all implemented and configured as
independent instantiations following the patterns described in Section
\ref{sec:UG:DirectInstantiation}.

In many cases, the simulation time is used as a fault's trigger variable and
independent variable.

\section{ Requirement-to-Test Mapping}

\par Here is an index of features and where they are tested:
\begin{flushleft}
\begin{enumerate}
 \item \textbf{Fault Effects}
  \begin{enumerate}
   \item \textbf{Profile Types}
    \begin{enumerate}

     \item \textbf{Holding}
      \begin{itemize}
       \item \textit{SIM\_fault\_manager : RUN\_overwrite}
         (Section \ref{sec:verif:FM:overwrite})
       \item \textit{SIM\_fault\_manager : RUN\_triggers}
         (Section \ref{sec:verif:FM:triggers})
      \end{itemize}

     \item \textbf{Fixed bias}
      \begin{itemize}
       \item \textit{SIM\_fault\_manager : RUN\_bias}
         (Section \ref{sec:verif:FM:bias})
       \item \textit{SIM\_fault\_manager : RUN\_function}
         (Section \ref{sec:verif:FM:function})
      \end{itemize}

     \item \textbf{Fixed-scale}
      \begin{itemize}
       \item \textit{SIM\_fault\_manager : RUN\_scale}
         (Section \ref{sec:verif:FM:scale})
      \end{itemize}

     \item \textbf{Stale-data}
      \begin{itemize}
       \item \textit{SIM\_fault\_manager : RUN\_stale}
         (Section \ref{sec:verif:FM:stale})
      \end{itemize}

     \item \textbf{Noise}
      \begin{itemize}
       \item \textit{SIM\_fault\_manager : RUN\_white\_noise}
         (Section \ref{sec:verif:FM:white_noise})
       \item \textit{SIM\_fault : RUN\_noise\_*}
         (Section \ref{sec:verif:DI:noise})
      \end{itemize}

     \item \textbf{Walk-off}
      \begin{itemize}
       \item \textit{SIM\_fault\_manager : RUN\_function}
         (Section \ref{sec:verif:FM:function:linear}).
       \item \textit{SIM\_fault : RUN\_linear}
         (Section \ref{sec:verif:DI:function:linear})
      \end{itemize}

     \item \textbf{Random-walk}
      \begin{itemize}
       \item \textit{SIM\_fault\_manager : RUN\_random\_walk}
         (Section \ref{sec:verif:FM:randomwalk})
      \end{itemize}

     \item \textbf{Periodic-fluctuation}
      \begin{itemize}
       \item \textit{SIM\_fault\_manager : RUN\_function}
         (Section \ref{sec:verif:FM:function})
      \end{itemize}

     \item \textbf{Initial Fault}
      \begin{itemize}
       \item \textit{SIM\_fault\_manager : RUN\_initial\_bias}
         (Section \ref{sec:verif:FM:initial_fault})
      \end{itemize}
    \end{enumerate}
   \end{enumerate}


   \item \textbf{Periodic Profiles}
    \begin{enumerate}
     \item \textbf{Shaped profiles}
      \begin{itemize}
       \item \textit{SIM\_fault\_manager : RUN\_function}
         (Section \ref{sec:verif:FM:function})
      \end{itemize}

     \item \textbf{Independent data types}
      \begin{itemize}
       \item \textit{SIM\_fault\_manager : RUN\_function2}
         (Section \ref{sec:verif:FM:function2})
       \item \textit{SIM\_fault : RUN\_linear\_int})
         (Section \ref{sec:verif:DI:function:linear})
      \end{itemize}

     \item \textbf{Configurable Parameters}
      \begin{itemize}
       \item \textit{SIM\_fault\_manager : RUN\_function}
        (Section \ref{sec:verif:FM:function})
       \item \textit{SIM\_fault\_manager : RUN\_function2}
        (Section \ref{sec:verif:FM:function2})
       \item \textit{SIM\_fault : RUN\_linear\_int}
         (Section \ref{sec:verif:DI:function:linear})
      \end{itemize}

     \item \textbf{Function Parameters}
      \begin{itemize}
       \item \textit{SIM\_fault : RUN\_sinewave}
         (Section \ref{sec:verif:DI:sinewave})
      \end{itemize}
    \end{enumerate}


   \item \textbf{Compounding Faults}
    \begin{itemize}
     \item \textit{SIM\_fault\_manager : RUN\_function}
       (Section \ref{sec:verif:FM:function:superimposed},
       using \verb'var_double')
     \item \textit{SIM\_fault\_manager : RUN\_function2}
       (Section \ref{sec:verif:FM:function2}, using \verb'var_double')
    \end{itemize}

   \item \textbf{Faults On and Off}
    \begin{itemize}
     \item using \textit{set\_fault\_enabled}:
      \begin{itemize}
       \item \textit{SIM\_fault\_manager : RUN\_bias}
         (Section \ref{sec:verif:FM:bias})
       \item \textit{SIM\_fault\_manager : RUN\_stale}
         (Section \ref{sec:verif:FM:stale}).
      \end{itemize}

     \item using \textit{disable()}:
      \begin{itemize}
       \item \textit{SIM\_fault : RUN\_disable\_*\_indirect}
         (Section \ref{sec:verif:DI:disable})
      \end{itemize}

     \item using \textit{enabled} flag directly:
      \begin{itemize}
       \item \textit{SIM\_fault : RUN\_disable\_*}
         (Section \ref{sec:verif:DI:disable})
      \end{itemize}
    \end{itemize}

   \item \textbf{Monte Carlo Support}
    \begin{itemize}
     \item \textit{SIM\_fault\_manager : RUN\_white\_noise}
       (Section \ref{sec:verif:FM:white_noise})
     \item \textit{SIM\_fault\_manager : RUN\_random\_walk}
       (Section \ref{sec:verif:FM:randomwalk})
     \item \textit{SIM\_fault\_manager : ERROR\_no\_random\_seed\_specified}
       (Section \ref{sec:verif:FM:randomwalk})
    \end{itemize}

   \item \textbf{Modifying Faults}
    \begin{itemize}
     \item \textit{SIM\_fault : RUN\_sinewave}
       (Section \ref{sec:verif:DI:sinewave})
     \item \textit{SIM\_fault : RUN\_new\_rate}
       (Section \ref{sec:verif:DI:new_rate})
     \item \textit{SIM\_fault\_manager} has examples in multiple runs.
    \end{itemize}
  \end{enumerate}

 \item \textbf{Fault Causes and Triggers}
  \begin{enumerate}
   \item \textbf{Comparators}
    \begin{itemize}
     \item \textit{SIM\_fault\_manager : RUN\_triggers}
       (Section \ref{sec:verif:FM:triggers})
    \end{itemize}

   \item \textbf{Compound Logical Comparisons}
    \begin{itemize}
     \item \textit{SIM\_fault\_manager : RUN\_stale}
       (Section \ref{sec:verif:FM:stale:multigroup})
     \item \textit{SIM\_fault\_manager : RUN\_random\_walk}
       (Section \ref{sec:verif:FM:randomwalk} using \verb'var_double')
    \end{itemize}

   \item \textbf{Triggering Variables}
    \begin{enumerate}
     \item \textbf{Data Types}
      \begin{itemize}
       \item \textit{SIM\_fault\_manager : RUN\_triggers}
         (Section \ref{sec:verif:FM:triggers})
       \item \textit{SIM\_fault : RUN\_string\_trigger}
         (Section \ref{sec:verif:DI:string_trigger})
      \end{itemize}
     \item \textbf{Configure Trigger-Values}
      \begin{itemize}
       \item using \textit{set\_trigger\_value}
        \begin{itemize}
         \item \textit{SIM\_fault\_manager : RUN\_triggers}
           (Section \ref{sec:verif:FM:triggers})
         \item \textit{SIM\_fault\_manager : RUN\_bias}
           (Section \ref{sec:verif:FM:bias})
        \end{itemize}
       \item setting value directly:
        \begin{itemize}
         \item \textit{SIM\_fault : all runs} (Section \ref{sec:verif:DI})
        \end{itemize}
      \end{itemize}
    \end{enumerate}

   \item \textbf{Trigger Controls}
    \begin{enumerate}
     \item \textbf{Enabling / Disabling triggers}
      \begin{itemize}
       \item \textit{SIM\_fault\_manager : RUN\_bias}
        (Section \ref{sec:verif:FM:bias})
      \end{itemize}

     \item \textbf{Limited Firing}
      \begin{itemize}
       \item \textit{SIM\_fault\_manager : RUN\_triggers}
         (Section \ref{sec:verif:FM:triggers})
       \item \textit{SIM\_fault\_manager : RUN\_bias}
         (Section \ref{sec:verif:FM:bias})
       \item \textit{SIM\_fault\_manager : RUN\_stale}
         (Section \ref{sec:verif:FM:stale})
      \end{itemize}
     \item \textbf{Periodic Triggers}
      \begin{itemize}
       \item \textit{SIM\_fault : RUN\_linear\_periodic}
         (Section \ref{sec:verif:DI:function:linear})
       \item \textit{SIM\_fault\_manager : RUN\_bias}
         (Section \ref{sec:verif:FM:bias})
      \end{itemize}
     \item \textbf{Multiple Firings}
      \begin{itemize}
      \item \textit{SIM\_fault\_manager : RUN\_bias}
       (Section \ref{sec:verif:FM:bias})
      \item \textit{SIM\_fault\_manager : RUN\_stale}
       (Section \ref{sec:verif:FM:stale})
      \item \textit{SIM\_fault\_manager : RUN\_random\_walk}
       (Section \ref{sec:verif:FM:randomwalk})
      \end{itemize}
    \end{enumerate} % end of Trigger Controls
  \end {enumerate} % end of Fault Causes


 \item \textbf{Fault Application}. Confirmation that the target-variable can
   be any of the specified data types is handled in all runs within
   \textit{SIM\_fault\_manager}. The manager is configured to provide a fault
   for each data-type of target variable.

 \item \textbf{Fault Implementation}
  \begin{enumerate}
   \item \textbf{Model-based Implementation}
     The encapsulation of faults within a class is seen in the
     \textit{FaultArchitectureTestSimObject} class in both
     \textit{SIM\_fault} (Section \ref{sec:verif:DI}), and
     \textit{SIM\_fault\_manager} (Section \ref{sec:verif:FM}).
     These examples are structured to emulate the way in which faults
     might be implemented in a model class.

   \item \textbf{Simulation-based Implementation} Both test simulations also
     demonstrate that the model can be implemented within a simulation.
     The only significant difference between a model-based implementation
     and implementation at the top-level of a simulation is access to data.
     In \textit{SIM\_fault\_manager : RUN\_triggers} (Section
     \ref{sec:verif:FM:triggers}), the faults are configured to
     act on private and const variables, and can adjust the values of
     variables with these protections.

   \item \textbf{Independent Faults} The structure found in \textit{SIM\_fault}
     (Section \ref{sec:verif:DI}) uses the architecture of independently
     instantiated, configured, and executed faults.

   \item \textbf{Managed Faults} The structure found in
     \textit{SIM\_fault\_manager} (Section \ref{sec:verif:FM})
     uses the XML-file-based definition and the fault-manager architecture to
     instantiate, configure, and execute faults.
     Of particular note, \textit{RUN\_function} tests the application of faults
     using the four identifying locations for regular updates:
     \begin{itemize}
      \item UPSTREAM
      \item INTERMEDIATE\_1
      \item INTERMEDIATE\_2
      \item DOWNSTREAM
     \end{itemize}
     and \textit{RUN\_initial\_bias} test the application of a fault using the
     identifying location for a Fault to be applied at initialization.
     \begin{itemize}
      \item INIT
     \end{itemize}
 \end{enumerate}

Additional model features are mapped out here:
\begin{itemize}
 \item Making trigger values randomized:
  \begin{itemize}
   \item \textit{SIM\_fault\_manager : RUN\_bias}
     (Section \ref{sec:verif:FM:bias})
  \end{itemize}
 \item Handling of invalid XML files or invalid entries in the XML file:
  \begin{itemize}
   \item \textit{SIM\_fault\_manager : RUN\_errors}
     (Section \ref{verif_errors})
   \item \textit{SIM\_fault\_manager : FAIL\_invalid\_xml}
      (Section \ref{verif_fails})
   \item \textit{SIM\_fault\_manager : FAIL\_bad\_filename}
     (Section \ref{verif_fails})
  \end{itemize}
\end{itemize}
\end{flushleft}

\section{Code Coverage}

Code coverage is illustrated in Figure \ref{fig:coverage}
\begin{figure}[ht]
  \centering
  \includegraphics[scale=0.45]{graphics/code_coverage.png}
  \caption{Code Coverage}
  \label{fig:coverage}
\end{figure}

\subsection{Exceptions}
\begin{itemize}
 \item \textbf{fault\_function.hh} (3 lines). False flag with the situation
       in which the function type is undefined (switch case \textit{default}).
       \newline These lines are all exercised in
       \textit{SIM\_fault/ERROR\_invalid\_function\_type}.
       Execution of this run as part of unit-testing has been confirmed.
       The cause of this false flag is unknown.
 \item \textbf{independent\_variable.hh} (1 line). False flag with the situation
       in the the independent variable is initialized twice (line 59).
       The test case \textit{SIM\_fault/ERROR\_double\_initialization} is
       specifically configured to hit this line.
       Execution of this run as part of unit-testing has been confirmed.
       The cause of this false flag is unknown.
 \item \textbf{trigger.hh} (5 lines)
 \begin{itemize}
   \item Destructors (3 lines).
         Standard omission of destructor execution due to testing process.
   \item False flag (1 line) in \textit{bool compare()} (line 183) with the
         default case of the Operator switch).
         This line is exercised in
         \textit{SIM\_fault/RUN\_string\_trigger\_GT}; this may
         have been optimized out because it has no effect, setting
         \textit{success} to the same value it was previously assigned.
   \item False flag in \textit{bool Trigger<T>::compare()} (line 240) with
         the default case of the Operator switch). This line
         is exercised in \textit{SIM\_fault/RUN\_linear\_invalid\_comparator};
         this may have been optimized out for the same reason as the previous
         false flag.
 \end{itemize}
 \item \textbf{fault\_manager.cc} (9 lines):
 \begin{itemize}
   \item Destructor (8 lines)
       Standard omission of recording the execution of destructors,
       this is a known issue with the testing process.  This method has been
       manually verified in gdb.
   \item Sanity check (1 line) in \textit{FaultManager::parse()} (line 281).
         This is
         an unreachable sanity-check error message, used to detect an error
         whereby an XML file is specified but has no XML content. It is left in
         place in case there are unknown methods for corrupting the XML file
         that could cause a failure here.
         It has been manually tested in gdb.
 \end{itemize}
 \item \textbf{fault\_bias.hh, fault\_overwrite.hh, fault\_scale.hh,
       fault\_stale.hh} (1 line each).
       Standard omission of recording the execution of (empty) destructors,
       this is a known issue with the testing process, but these methods are all
       empty.
\end{itemize}



\section{SIM\_fault\_manager} \label{sec:verif:FM}
\par RUN\_fault\_manager tests each type of fault on all the valid variable
types listed in Section \ref{sec:intro:faults}:
\begin{codeC}
char               var_char;
unsigned char      var_uchar;
short              var_sint;
unsigned short     var_usint;
int                var_int;
unsigned int       var_uint;
long               var_lint;
unsigned long      var_ulint;
long long          var_llint;
unsigned long long var_ullint;
float              var_float;
double             var_double;
bool               var_bool;
\end{codeC}

The fault-types are investigated one at a time in the associated runs:
\begin{enumerate}
\item RUN\_bias
\item RUN\_scale
\item RUN\_overwrite
\item RUN\_stale
\item RUN\_white\_noise
\item RUN\_random\_walk
\item RUN\_function
\end{enumerate}

The following runs complete the testing for this run, testing general model
functionality rather than specific fault-types.
\begin{enumerate}[resume]
\item RUN\_function2
\item RUN\_initial\_bias
\item RUN\_errors
\item RUN\_triggers
\item RUN\_trigger\_limited
\end{enumerate}

In most runs, the variables are each assigned a constant value every timestep.
Unless a variable is actively being faulted, it will retain that value, and
deviations from it indicate application of a fault.
In \textit{RUN\_stale}, \textit{RUN\_initial\_bias}, and \textit{RUN\_triggers},
we make an exception to this pattern
because these runs require the non-faulted values to vary by known increments
in order to demonstrate the success of the fault.

For logging purposes, the variable \verb"var_char" is cast to an
\textit{int} variable called \verb"var_char_numeric" to ensure
it is logged as a number instead of an ASCII character.

\subsection{RUN\_bias} \label{sec:verif:FM:bias}
\par In RUN\_bias, a bias fault is created for each of the valid types. The
faults are defined in the file \textit{bias.xml}.

\subsubsection{ Regular Triggers}
Most faults use the trigger called \textit{bias trigger}, which is set to
trigger when $t=1$. These faults should therefore show a bias being
applied at $t=1$, and holding thereafter.
In the input file, this trigger has its
trigger-value reset to 4.0 at $t=3.0$. All faults using this trigger should
therefore return to their non-triggered state at $t=3.0$, re-trigger at $t=4.0$,
and hold this state to the end of the run.
\verb"var_char", \verb"var_uchar", \verb"var_usint", \verb"var_llint", and
\verb"var_double" show this behaviour without any modification.
Faults that use other modifications and triggers are discussed below.
Results of these tests are shown in Figure \ref{fig:verif:FM:bias}

\subsubsection {Modified Parameters}
The \textit{set\_param} method is exercised via the manager's
\textit{set\_fault\_param} method.
First, the method is given an invalid parameter name, which causes an error to
be printed to the console.

\verb"var_sint" uses trigger \textit{bias trigger} as discussed above, but at
$t=5$, the \textit{set\_fault\_param} method is used to modify the \textit{bias}
parameter from 150 to 100. \verb"var_sint" should show the same on/off pattern
as seen in the other variables using \textit{bias trigger}, but the output value
should show a step function at $t=5$ as the modifier is adjusted from $150$
to $100$.

\subsubsection{Disable Fault}
\verb"var_float"'s fault uses trigger \textit{bias trigger} as discussed above;
this fault is disabled at the beginning of the run and enabled at $t=5$. This
fault is applied at $t=5$, with the triggering criterion also being satisfied at
$t=5$.

\subsubsection{Disable Trigger}
\verb"var_uint"'s fault uses trigger \textit{bias trigger} as discussed above;
at the beginning of the run, the \textit{bias trigger} trigger is disabled in
this fault with the \verb'set_fault_trigger_enabled' function. Note that this
does not affect the operation of \textit{bias trigger} in any other fault. At
$t=5$, this trigger is re-enabled using the same function. Disabling the fault's
only trigger should have the same effect as disabling the fault in the
previous example; this fault should similarly trigger only at $t=5$.


\subsubsection {Fire Limit}
\verb"var_int"'s fault uses trigger \textit{bias trigger} as discussed above,
but is given a fire limit of 9.
Because the simulation updates at 10 Hz, this fault should only be triggered for
0.9 seconds.

\subsubsection {Periodic Triggers}
\par The faults corresponding to \verb"var_lint" and \verb"var_ulint" use
periodic triggers.
The trigger for \verb"var_lint" has a trigger-value of 1, a period of 3, and a
length of 1.5. It should trigger at $t=1$ and remain triggered until
$t=1+1.5=2.5$. It should re-trigger at $t=1+3=4$ and remain triggered until
$t=4+1.5=5.5$.

The fault for \verb"var_ulint" demonstrates how periodic triggers behave when
the trigger variable has a negative value, and follows a decreasing function.
This fault uses the trigger called \textit{negative periodic trigger}, which
uses \verb"negative_time" as its trigger variable, and triggers when this
variable is $\leq -1$. \verb"negative_time" is simply the
negative value of the simulation time. As before, this fault uses values for
period of $3$, and length of $1.5$.
Because periodic triggers use the absolute-value of \verb"std::fmod" to evaluate
the fractional position within a period, the sign change has no effect and
\verb"var_ulint" should behave the same way as \verb"var_lint".
For $t_- =-(3n + \Delta)$, we have the application of \verb'fmod' yielding
$|t_- \% 3| = |-\Delta| = \Delta$, the same result from applying the modulo function
to the positive time: $|t_+ \%3| = |(3n + \Delta) \%3| = \Delta$

\subsubsection {Random Trigger}
\par The fault corresponding to \verb"var_ullint" has a trigger with a random
trigger value.
The distribution is flat with a minimum of 1.5 and maximum of 5.5, so the fault
should be triggered somewhere between those two times. Note -- for
reproducibility we use a fixed seed, which means for this test the fault always
triggers at the same time; but the output has been manually reviewed using
different seeds to confirm that the trigger-time does fluctuate.

\par The results of RUN\_bias are shown in Figure \ref{fig:verif:FM:bias}.
All the variables show the expected behaviour described in this section.

\begin{figure}
  \centering
  \includegraphics[scale=0.5]{graphics/verif_bias.png}
  \caption{RUN\_bias Results}
  \label{fig:verif:FM:bias}
\end{figure}
\clearpage

\subsection{RUN\_scale} \label{sec:verif:FM:scale}
\par In RUN\_scale, a scale fault is created for each of the valid types.
These faults all use the same trigger, which has a value of 1.

\par The \textit{set\_param} method is exercised via the manager's
\textit{set\_fault\_param} method.
First, the method is given an invalid parameter name, which causes an error to
be printed to the console.
Later, at $t=2$, it is used to change \verb"var_sint"'s fault's scale factor
from 150 to 100.

\par The results of RUN\_scale are shown in Figure \ref{fig:scale}.
All the variables show the expected behaviour.
(Note that \verb"var_char" has an unfaulted value of 97 and its fault has a
scale factor of 2.
\verb"char"s can be between -128 and 127, so the faulted value overflows and is
interpreted as -62 instead of 194.)

\begin{figure}[hb]
  \centering
  \includegraphics[scale = 0.58]{graphics/verif_scale.png}
  \caption{RUN\_scale Results}
  \label{fig:scale}
\end{figure}
\clearpage

\subsection{RUN\_overwrite} \label{sec:verif:FM:overwrite}
\par In RUN\_overwrite, an overwrite fault is created for each of the valid
types.
These faults all use the same trigger, which triggers at $t=1$.

For most faults, the overwrite parameter is configured in the XML file
\textit{xml\_files/overwrite.xml} to have a value of 150. The exceptions are:
\begin{itemize}
 \item \verb'var_char' which has an overwrite value of -62.
 \item \verb'var_uchar' which has an overwrite value of 96.
 \item \verb'var_float' which takes a random value with mean of 150 and
 standard-deivation 10; with the specified seed, this generates a
 pseudo-random value of $144.50253$
 \item \verb"var_ullint" starts at 150 then at $t=2$ the
 \textit{set\_fault\_param} method is called from the input file to reset this
 parameter to 40.
 \item \verb"var_bool" which has an override value of 1 (true).
\end{itemize}

Within this run, we also exercise the \textit{set\_fault\_param} method
with invalid parameter name. This causes the printing of an error
to the console.


\par The results of \textit{RUN\_overwrite} are shown in Figure
\ref{fig:verif:FM:overwrite}.
All the variables show the expected behaviour.

\begin{figure}[hb]
  \centering
  \includegraphics[scale=0.5]{graphics/verif_overwrite.png}
  \caption{RUN\_overwrite Results}
  \label{fig:verif:FM:overwrite}
\end{figure}
\clearpage

\subsection{RUN\_stale} \label{sec:verif:FM:stale}
\par In RUN\_stale, a stale-value fault is created for each of the valid types.
These faults all use the same trigger, which triggers at $t=1$.
When a stale-value fault is triggered, it maintains its variable at a constant
value.
To observe this effect, this run nominally changes each of the variables
\verb'var_*' at the beginning of each timestep: the numeric values are
incremented by $1$ and the Boolean variable \verb'var_bool' is flipped each
cycle. When the stale fault is applied, these changes will be overridden by the
fault.

\subsubsection{Single Trigger}
The fault corresponding to most variables use a single trigger. For most
variables, their values will increase until $t=1.0$, then hold constant
after its fault has triggered.

Note that \verb'var_bool' will oscillate between $0$ and $1$ for $t<1$
(rather than increasing with the other variables) and
hold constant thereafter.

\subsubsection{Multiple triggers}
\par The fault corresponding to \verb"var_uchar" is given an additional
trigger in the same trigger-group; this second trigger is triggered for $t
\leq 2.0$. Because all triggers in a group must be triggered for the group
to be triggered, this trigger-group reverts to non-triggered  when the
second trigger reverts to non-triggered at $t>2$, resulting in the fault
being removed at that time.

\subsubsection{Multiple trigger-groups} \label{sec:verif:FM:stale:multigroup}
\verb'var_sint' is configured with two trigger-groups; the first group is
identical to that used in the Multiple Triggers test (\verb"var_uchar"),
as outlined above.
The second trigger-group adds a new trigger with constraint $t\geq3$.
As with the Multiple Trigger test (\verb"var_uchar"), the fault is
triggered at $t=1$ and reverted at $t=2$ as the first trigger-group is triggered
and reverted. With the triggering of the second trigger-group at $t=3$, the
fault is re-triggered because a fault is triggered when \textit{any} of its
trigger-groups are triggered.

\subsubsection{ Disable and Re-enable}
The fault corresponding to \verb"var_usint" uses a single trigger (with
characteristics matching those of the single-trigger identified above) but it
is disabled at $t=2$ and re-enabled at $t=3$ using the \verb'set_fault_enabled'
function. Note that this
function call when disabling a fault forwards to the fault's \textit{disable()}
method, which results in the fault resetting when it is re-enabled. This is one
step more than simply setting \textit{enabled=false}, which would not result in
a reset but requires direct access to the fault.

Note also the potentially unexpected visual at $t=2$: the fault is disabled
at the very top of the frame; within the frame, the value of \verb"var_usint"
is incremented, then the fault is tested and applied. At $t=2.0$, the value of
\verb"var_usint" is incremented from its value at $t=1.9$, then this
incremented value is not reset to its faulted
value because the fault has already been disabled. Consequently, the last
faulted data point is at $t=1.9$, and by $t=2.0$ the variable is already in
its increasing phase; the `corner' of the plotted evolution of
\verb"var_usint" is seen at $t=1.9$, not $t=2.0$.
Re-enabling the fault at $t=3$ results in the fault-value being
recorded when it is first processed, at $t=3.0$. Consequently,
\verb"var_usint" after $t=3.0$ retains the value from $t=3.0$; the
`corner' of the plotted evolution of \verb"var_usint" as it is re-enabled is
seen at $t=3.0$.

\subsubsection {Fire Limit}
The fault applied to \verb"var_int" also uses the single trigger outlined
above, with the additional configuration of being given a fire
limit of 7. In a simulation updating at 10 Hz, this translates to 0.7 seconds,
resulting in the fault disabling itself at $t=1.7$. The fault is re-enabled in
the input file at $t=2.5$.
It should disable itself again at $t=3.2$.
The same `corner` explanation from the previous test applies here, with the
corners in the data appearing at $t=1.0, 1.6, 2.5, 3.1$.

\subsubsection {Parameter Setting}
\par The \textit{set\_param} method is exercised via the manager's
\textit{set\_fault\_param} method.
Stale-value faults have no settable parameters, so the method simply prints an
error to the console.

The results of RUN\_stale are shown in Figure \ref{fig:verif:FM:stale}.
All the variable show the expected behaviour described in this section.

\begin{figure}[hb]
  \centering
  \includegraphics[scale=0.58]{graphics/verif_stale.png}
  \caption{RUN\_stale Results}
  \label{fig:verif:FM:stale}
\end{figure}
\clearpage

\subsection{RUN\_white\_noise} \label{sec:verif:FM:white_noise}
In RUN\_white\_noise, white-noise faults are created for \verb"var_float"
and \verb"var_double".
These faults both use a trigger with a threshold \textit{value} of 0,
so they are triggered the entire run.

The fault for \verb"var_float" uses a Gaussian distribution.
The noise has a mean of $0.0$ and a standard deviation of $0.25$ in the
first half of
the run and a mean of $0.3$ and standard deviation of $0.02$ in the second
half.  The unfaulted value of \verb"var_float" is 0.1 so the faulted
variable value should have a mean of $0.1$ during the first half, and $0.4$
during the second half.

The fault for \verb"var_double" uses uniform distribution.
The noise has a range of $[-1, 1]$ in the first half of the run and
$[0.45, 0.55]$ in the second half. The first range is specified as a
(\textit{rel\_min, mean, rel\_max}) set.
The unfaulted value of \verb"var_double" is 0.2 so the faulted
variable value should fall between $[-0.8, 1.2]$ during the first half, and
$[0.65, 0.75]$ during the second half.

\textit{set\_param} (via \textit{set\_fault\_param}) is used to change the
distribution parameters halfway through the run.
An additional call to \textit{set\_param} uses an invalid parameter name,
which causes an error to be printed to the console.

For further testing, a set of 10 runs were collected, with each run using a
different seed for the random number generators. This exercise demonstrates
the feasibility of using the Monte Carlo engine to distribute the
white=-noise faults; the setting of the seed in the input file (or
monte-input file) takes precedence over the setting of the seed in the XML
file.

The results of RUN\_white\_noise are shown in Figures
\ref{fig:verif:FM:whitenoise}, \ref{fig:verif:FM:whitenoise2}, and
\ref{fig:verif:FM:whitenoise3}.
Both variables behave as expected.

\begin{figure}[htb]
  \centering
  \includegraphics[scale=0.5]{graphics/verif_white_noise.png}
  \caption{Single-run RUN\_white\_noise Results}
  \label{fig:verif:FM:whitenoise}
\end{figure}

\begin{figure}[htb]
  \centering
  \includegraphics[scale=0.5]{graphics/verif_white_noise_MC.png}
  \caption{10-run RUN\_white\_noise Results}
  \label{fig:verif:FM:whitenoise2}
\end{figure}

\begin{figure}[htb]
  \centering
  \includegraphics[scale=0.5]{graphics/verif_white_noise_MC2.png}
  \caption{10-run RUN\_white\_noise Results}
  \label{fig:verif:FM:whitenoise3}
\end{figure}
\clearpage

\subsection{RUN\_random\_walk} \label{sec:verif:FM:randomwalk}
\par In RUN\_random\_walk, random-walk faults are created for \verb"var_float"
and \verb"var_double".
The triggers are set up such that \verb"var_float"'s fault is always triggered
and \verb"var_double"'s is always triggered except when $8<t<9$.
The random variables used in the random-walk function (see Section
\ref{sec:intro:randomwalk}) have a Gaussian distribution in
\verb"var_float"'s fault
and a uniform distribution in \verb"var_double"'s fault.

The random variable in \verb"var_float"'s fault starts with a mean of 0 and
a standard deviation of 1.
At $t=6$, the mean is changed to 0.5 and the standard deviation to 0.
After this point, \verb"var_float" should increase linearly by 0.5 every
cycle (at 10Hz, that equates to an increase of 5 per second).

The random variable in \verb"var_double"'s fault starts with a range of
$[-1.732, 1.732]$.
This is approximately $\pm\sqrt{3}$, which makes the standard deviation
approximately 1.
At $t=6$, the range is changed to $[0.45, 0.55]$; after this point,
\verb"var_double" should increase by approximately 5 per second.
When the fault is no longer triggered, it should reset the random walk to 0 and
start from that point when it is re-triggered.

\textit{set\_param} (via \textit{set\_fault\_param}) is used to change the
distribution parameters at $t=6$.
An additional call to \textit{set\_param} uses an invalid parameter name,
which causes an error to be printed to the console.

For further testing, a set of 10 runs were collected, with each run using a
different seed for the random number generators. This exercise demonstrates
the feasibility of using the Monte Carlo engine to distribute the
random-walk faults; the setting of the seed in the input file (or
monte-input file) takes precedence over the setting of the seed in the XML
file.

The results of RUN\_random\_walk are shown in Figures
\ref{fig:verif:FM:randomwalk}, \ref{fig:verif:FM:randomwalk2}, and
\ref{fig:verif:FM:randomwalk3}.
Both variables behave as expected.

\begin{figure}[ht]
  \centering
  \includegraphics[scale=0.5]{graphics/verif_random_walk.png}
  \caption{Single-run RUN\_random\_walk Results}
  \label{fig:verif:FM:randomwalk}
\end{figure}

\begin{figure}[ht]
  \centering
  \includegraphics[scale=0.5]{graphics/verif_random_walk_MC.png}
  \caption{10-run RUN\_random\_walk Results}
  \label{fig:verif:FM:randomwalk2}
\end{figure}

\begin{figure}[ht]
  \centering
  \includegraphics[scale=0.5]{graphics/verif_random_walk_MC2.png}
  \caption{10-run RUN\_random\_walk Results}
  \label{fig:verif:FM:randomwalk3}
\end{figure}

The run \textit{ERROR\_no\_random\_seed\_specified} is a derivative of
\textit{RUN\_random\_walk}, for which the seed is specified to $0$.
Setting the seed to $0$ is functionally equivalent to not setting it at all
because the default value is also $0$.
With $seed=0$, the model will generate a seed from the
system clock. This is not truly an error test-case, but it is labeled as such
because it is not suitable for regression testing --
the output from this run changes with every execution, as expected.
\clearpage


\subsection{RUN\_function} \label{sec:verif:FM:function}
\par In RUN\_function, a function fault is created for each of the valid types.

One of two independent variables is used for each function:
\begin {itemize}
  \item \textit{current\_time} $(t)$ has a domain $[0,70]$
  \item \textit{offset\_time} $t'=t-30$ is offset from \textit{current\_time}
    and has domain $[-30,40]$. This supports testing of the sign-transition in
    the independent variable (when its value is configured to be
    \textit{absolute}, if it is configured to be \textit{relative}, it will
    start from $0$ and be positive thereafter).
\end{itemize}

This run utilizes 12 of the 13 variables (excluding \verb'var_bool')
to cover a broad distribution of the
available function options. A description of the mapping between function
configurations and specific variables follows. The results of all 12
variables are shown in Figures \ref{fig:function1} and \ref{fig:function2}
following the description.


\begin{enumerate}
\subsubsection{Linear Function} \label{sec:verif:FM:function:linear}
 \item \verb'var_char' (logged as \verb'var_char_numeric') uses the trigger
   simply named \textit{trigger}, which triggers when $t\geq5$. This fault uses:

   \begin {itemize}
     \item \textit{current\_time} $(t)$ as its independent variable, with
       parameters:
     \item nominal value 0
     \item rate 10
   \end{itemize}
   The output value starts to ramp up at $t=5$; the variable overflows at $127$
   and wraps back to $-128$, producing a saw-tooth profile.

 \item \verb'var_uchar' uses the same configuration as \verb'var_char', except
   the rate is set to $-10$ and the \textit{nominal} keyword is switched for the
   \textit{initial} keyword in setting the parameters (these two keywords
   are interchangeable). It produces a similar
   profile, except the saw-tooth
   pattern is reversed by the sign of the rate, and the unsigned nature of the
   variable shifts the range to non-negative values, $[0,255]$.

 \item \verb'var_sint' uses the same configuration as \verb'var_char', except
   the nominal value is set to $200$. This fault introduces a step at $t=5$ and
   smooth slope thereafter.

 \item \verb'var_usint' is also based on the configuration of \verb'var_char',
   but uses \textit{test\_object.offset\_time} as the independent
   variable for the linear-function, set to be used in a \textit{relative}
   configuration. In its relative configuration, this will be interpreted as
   starting from $0$ and produce a smooth transition to the linear ramp.
   Additionally, for this configuration, neither \textit{nominal} nor
   \textit{initial} keywords were used to specify the zero value of the
   y-intercept. This value defaults to zero without specification.

 \item \verb'var_int' also uses \textit{test\_object.offset\_time} as the
   independent variable but in an absolute configuration. At $t=5$,
   $t'-5-30=-25$ and the fault function will have a value $f=0+10(-25)=-250$.
   The output should step down, then transition to its smooth up-ramp.

\subsubsection{Periodic Functions}
Note -- for periodic functions, the difference between treating the independent
variable as \textit{absolute} or \textit{relative} is largely moot.
The function is evaluated
indirectly from the independent variable, with the independent variable used to
generate the phase of the periodic function, and the phase used to generate its
value. This phase
is always incremented (approximation to being numerically integrated) based on
the difference between the current and previous values of the independent
variable; additionally, the incremented/integrated component of the phase always
starts at $0.0$. In a sense then, the independent value is always treated as
being \textit{relative} when using periodic functions. Consequently, without
loss of generality, we can always use \textit{current\_time} (rather than
\textit{offset\_time} as the independent variable, and must apply an initial
phase offset to introduce a sign-change in the value of phase.

\subsubsection{Square Wave}
 \item \verb'var_uint' uses the standard \textit{trigger} with
   \textit{test\_object.current\_time} as the
   independent variable, in a relative configuration. The square-wave is
   configured to have an amplitude of $4$ and frequency of $0.1$ (i.e. period
   is $10$). At $t=5$, the square-wave should initiate and continue for the
   duration of the run, starting with a step up at phase $0.0$. Note on
   verification data, not seen in plots for lack of resolution: the steps taken
   at phase = $\phi = {0, 0.5}$ may be offset by one frame from expectation
   based on the rounding associated with numerical integration to these critical
   values. The result for
   $\phi=0.499999999$ and
   $\phi=0.500000001$ are well defined and distinct, making the result for
   $\phi=0.500000000$ ambiguous depending on which side of the line the
   approximated value falls.

\subsubsection{Absolute Independent Variable}
 \item \verb'var_lint' uses a square-wave fault configured similarly to the
   previous case, except that it uses \textit{test\_object.current\_time} as
   the independent variable in an absolute configuration. As discussed above
   (under \textit{Periodic Functions}), this should make no difference and the
   results from this fault should mirror those from \verb'var_uint' (albeit
   shifted up by 1 because the nominal value of \verb`var_lint` prior to
   application of the fault is 1 larger than that of \verb'var_uint').

\subsubsection{Offset Parameter / Zero-crossing}
 \item \verb'var_ulint' is similar to  \verb'var_uint' except that the phase
   is now provided a phase-offset of $-2.5$ (actually, $-2.49999999$ to avoid
   the ambiguity of the step function that is applied at a half-integer phase).
   This fault should replicate a similar square-wave pattern, this time
   verifying that the square-wave correctly computes as the phase changes sign -
   (it starts negative due to the phase-offset and continues to increase through
   zero).
   Note also that because the phase starts at $-2.49999999$ (or equivalently,
   at $+0.50000001$), the square-wave should start at its lower value, with a
   step-down.

\subsubsection{Evolving Parameters}
 \item \verb'var_llint' is similar to  \verb'var_ulint' except that the phase
   offset is now also given a rate of $0.1$, meaning that the phase-offset
   continually increases. This should increase the rate at which the phase
   advances, effectively doubling the frequency (the phase advances by 0.1 per
   second due to integrating the constant frequency, and by another 0.1 per
   second because the phase-offset is increasing by that rate).

\subsubsection{Evolving Parameters with Overwrites and Independent
               Independent-Variable}
 \item \verb'var_ullint' returns the phase-offset to a constant but applies a
   rate to the amplitude and frequency. The amplitude is progressively decreased
   from $400$ with a rate of $-5$. The frequency is also set to decrease from an
   initial value of $0.35$ with a rate of $+0.01$ using
   \textit{negative\_time} as the independent variable for frequency.
   At $t=40$, the sign of the frequency rate is reversed (using the
   optional flag, see Section \ref{sec:formulation:fault:params}).
   For $5<t<40$, the width of the square-wave pattern should broaden as the
   frequency drops towards zero; for $40<t<70$, the width of the
   square-wave pattern should narrow as the frequency increases in magnitude
   again as it becomes increasingly negative.
   Throughout the run (for $t>5$), the amplitude decreases.

   Note that this variable (\verb'var_ullint') is an integer, so to make this
   pattern appear less discrete, we needed to increase the amplitude of this
   square wave. But it is also an unsigned integer, and simply increasing
   the amplitude would have pushed the resulting value negative, wrapping
   around to a large positive value.
   To overcome this problem, we apply fire-limited overwrites (single-shot,
   limited with \textit{fire\_limit = 1} to both \verb'constant\_ullint'
   and \verb'var\_ullint' using the same trigger. Consequently, the study of
   this pattern
   required an additional 2 faults to prepare it.

  \subsubsection {Triangle Wave}
  \item \verb'var_float' is similar to \verb'var_ulint' except that the function
    pattern is switched to a triangle-wave.

  \subsubsection {Sine Wave}
  \item \verb'var_double' uses a similar configuration, but with a sine-wave.
  This simple pattern is retained until $t=35$

  \subsubsection {Superimposed Functions} \label{sec:verif:FM:function:superimposed}
  \verb'var_double' continues for $t>35$ with the application of a second
    fault, this time a square-wave with a period double that of the sine-wave.
    This applies a $\pm1$ shift to the sine-wave, alternating between subsequent
    periods of the sine-wave.

\end{enumerate}
\begin{figure}[ht]
  \centering
  \includegraphics[scale=0.55]{graphics/verif_function1.png}
  \caption{RUN\_function Results (Part I)}
  \label{fig:function1}
\end{figure}
\begin{figure}[ht]
  \centering
  \includegraphics[scale=0.55]{graphics/verif_function2.png}
  \caption{RUN\_function Results (Part II)}
  \label{fig:function2}
\end{figure}


The results of RUN\_function are shown in Figure \ref{fig:function1} and
\ref{fig:function2}.
All the variables show the expected behaviour described in this section.

\clearpage

\subsection{RUN\_function2} \label{sec:verif:FM:function2}
\textit{RUN\_function2} is a simpler set of tests than \textit{RUN\_function},
targeted at evaluating the template nature of the independent variables. As in
\textit{RUN\_function}, a set of linear function faults is established to apply
to the different variable types. In \textit{RUN\_function}, the independent
variable for the faults was simply the simulation time, allowing the focus
to be on evaluating the function behavior. For this set of tests, we
want to test the ability to use different independent variables; for each
fault-function, we assign the independent variable to be the output of one of
the faulted variables.

To keep the situation analytically verifiable, most linear functions are either
$y=2x$ or $y=0.5x$, resulting in the step-by-step changes being in powers of 2.

The configuration of the test is shown in Table \ref{tbl:function2} and the
results in Table \ref{tbl:function2_results}. Because this test is only
evaluating whether different variables can be used as independent variables,
it is only run for a few steps to confirm the expected pattern.

\begin{table}[ht]
\centering
\caption{RUN\_function2 Configuration} \label{tbl:function2}
\begin{tabular}{|l|c|c|c|} \hline
\textbf{variable} & \textbf{independent variable} & \textbf{rate} &
\textbf{step-difference} \\ \hline
var\_char & current\_time & 1 & 1 \\ \hline
var\_uchar & var\_char & 2 & 2 \\ \hline
var\_sint & var\_uchar & 2 & 4 \\ \hline
var\_usint & var\_sint & 2 & 8 \\ \hline
var\_int & var\_usint & 0.5 & 4 \\ \hline
var\_uint & var\_int & 0.5 & 2 \\ \hline
var\_lint & var\_uint & 0.5 & 1 \\ \hline
var\_ulint & var\_lint & 2 & 2 \\ \hline
var\_llint & var\_ulint & 2 & 4 \\ \hline
var\_ullint & var\_llint & 2 & 8 \\ \hline
var\_float & var\_ullint & 0.5 & 4 \\ \hline
var\_double & var\_float & 0.5 & 2 \\ \hline
var\_double & var\_double & 10 & 20 \\ \hline
var\_double & var\_bool & 2 & XXX \\ \hline
\end{tabular}
\end{table}
Note - \verb'var_double' is given 3 faults:
\begin {itemize}
 \item one using \verb'var_float'. This produces steps of 2.
 \item one using itself, \verb'var_double' to test tests whether a fault can be
       self- referential. With a rate of 10, this produces steps of 20 (10 times
       the 2 from the first fault)
 \item one using \verb'var_bool', a boolean variable. Boolean variables cannot be
       used as independent variables, this fault generates an error message.
\end {itemize}
The two valid faults combine to produce steps of 22 between each data point.


\begin{table}[ht]
\centering
\caption{RUN\_function2 Results} \label{tbl:function2_results}
\begin{tabular}{|l|c|c|c|c|c|c|c|} \hline
\textbf{current\_time}	 & \textbf{0.00} & \textbf{1.00} & \textbf{2.00} &
\textbf{3.00000} & \textbf{4.00000} & \textbf{5.00000} & Delta \\ \hline
var\_char	   & 98.00 & 99.00 & 100.00 & 101.00 & 102.00 & 103.00 & 1 \\ \hline
var\_uchar	 & 1.00  & 3.00  & 5.00   & 7.00   & 9.00   & 11.00  & 2 \\ \hline
var\_sint	   & 2.00  & 6.00  & 10.00  & 14.00  & 18.00  & 22.00  & 4 \\ \hline
var\_usint	 & 3.00  & 11.00 & 19.00  & 27.00  & 35.00  & 43.00  & 8 \\ \hline
var\_int	   & 5.00  & 9.00  & 13.00  & 17.00  & 21.00  & 25.00  & 4 \\ \hline
var\_uint	   & 5.00  & 7.00  & 9.00   & 11.00  & 13.00  & 15.00  & 2 \\ \hline
var\_lint	   & 8.00  & 9.00  & 10.00  & 11.00  & 12.00  & 13.00  & 1 \\ \hline
var\_ulint	 & 7.00  & 9.00  & 11.00  & 13.00  & 15.00  & 17.00  & 2 \\ \hline
var\_llint	 & 8.00  & 12.00 & 16.00  & 20.00  & 24.00  & 28.00  & 4 \\ \hline
var\_ullint	 & 9.00  & 17.00 & 25.00  & 33.00  & 41.00  & 49.00  & 8 \\ \hline
var\_float	 & 0.10  & 4.10  & 8.10   & 12.10  & 16.10  & 20.10  & 4 \\ \hline
var\_double	 & 0.20  & 22.20 & 44.20  & 66.20  & 88.20  & 110.20 & 22\\ \hline
\end{tabular}
\end{table}

\clearpage

\subsection{RUN\_initial\_bias} \label{sec:verif:FM:initial_fault}
In this run, we are testing the fault location \textit{INIT}, which should only
be applied during initialization. For this scenario, we use the file
\textit{xml\_files/initial\_bias.xml} to configure a Fault to
bias \verb'var_float' by $+2.5$ at initialization, with no other faults
registered. The scenario uses the natural incremental progression of the sim
variables rather than the typical reset-to-constant; this allows the initial
bias to be preserved. The scenario logs the affected variable \verb'var_float'
and \verb'var_double' as a control.
We should epect to see:
\begin{itemize}
\item \verb'var_float' initially incremented from its initial value of $0.1$ by
      $2.5$ from the bias; by the end of the first cycle, the incremental
      progression adds an additional $1.0$ to produce an initial logged value
      of $0.1+2.5+1.0 = 3.6$. Thereafter, the value should increase by $+1.0$
      per cycle, or $+10.0$ per second.
\item \verb'var_double' is used as the control. It starts at $0.2$ and is
      subject to the same natural progression, but without the initial bias.
      \verb'var_double' should therefore always log a value that is $2.4$ lower
      than \verb'var_float'.
\end{itemize}
These expectations are confirmed in Figure \ref{fig:verif:init_bias}
\begin{figure}[h!]
  \centering
  \includegraphics[scale=0.5]{graphics/verif_init_bias.png}
  \caption{RUN\_initial\_bias Results}
  \label{fig:verif:init_bias}
\end{figure}



\subsection{RUN\_triggers} \label{sec:verif:FM:triggers}
\par RUN\_triggers tests all the valid types of triggers (including string
triggers) and all of the possible operators.
The faults all act on the constant
values that are used to replace the variable values at the beginning of every
timestep. This confirms that faults can be applied to protected and const
variables. Each fault is also given a fire limit of 1; while the fault will
disable after a single application, there is no mechanism for changing the value
of these const variables so they remain in their faulted state.

The triggers are listed in Table \ref{tbl:triggers}.
\begin{table}[htb]
\centering
\caption{RUN\_trigger Faults} \label{tbl:triggers}
\begin{tabular}{|l|c|c|} \hline
\textbf{\#} & \textbf{Trigger Condition} & \textbf{Action} \\ \hline
1 & command $\neq$ "hold"  & constant\_bool = true \\ \hline
2 & var\_bool = true       & constant\_char = 100 \\ \hline
3 & var\_char > 99         & constant\_uchar = 5 \\ \hline
4 & var\_uchar = 5         & constant\_sint = 6 \\ \hline
5 & var\_sint $\neq$ 2     & constant\_usint = 0 \\ \hline
6 & var\_usint $\leq$ 1    & constant\_int = 1 \\ \hline
7 & var\_int < 2           & constant\_uint = 9 \\ \hline
8 & var\_uint $\geq$ 8     & constant\_lint = 10 \\ \hline
9 & var\_lint $\geq$ 10    & constant\_ulint = 11 \\ \hline
10 & var\_ulint $\geq$ 11  & constant\_llint = 12 \\ \hline
11 & var\_llint $\geq$ 12  & constant\_ullint = 13 \\ \hline
   &                       &                       \\ \hline
12 & var\_ullint $\geq$ 60 & constant\_float = 1.0001 \\ \hline
13 & var\_float $\neq$ 0.1 & constant\_double = 0.3 \\ \hline
14 & var\_double = 0.3     & constant\_bool = false \\ \hline
   &                       &                       \\ \hline
15 & command = "stop"      & constant\_uchar = 265 \\ \hline
\end{tabular}
\end{table}

\begin{itemize}
 \item \verb"command" is a std::string that is only used by this run; it is
   initialized to the value "hold". At $t=0.3$, it is changed to "go",
   triggering the first fault in Table \ref{tbl:triggers}. This fault sets
   \verb"constant_bool" to true.
 \item On the next timestep ($t=0.4$), the value of \verb"constant_bool" is
   copied to \verb"var_bool", leading to the triggering of the second fault
   which sets \verb"constant_char" to $100$.
 \item At $t=0.5$, \verb"constant_char" is copied to \verb"var_char" and the
   third fault is triggered.
 \item The faults continue to cascade, stopping at fault 11, which sets
   \verb"constant_ullint" to $13$, which is copied to \verb"var_ullint" at
   $t=1.4$. Fault 12 is monitoring for \verb"var_ullint" being $\geq 60$, which
   is not satisfied.
 \item at $t=1.7$, the trigger-value for the next value is reduced to $10.0$
   using the method \newline
   \textit{set\_trigger\_value(...)}. This continues the
   cascade of
   faults, triggering faults 13 and 14 in subsequent timesteps.
 \item At $t=1.9$, \verb"command" is assigned the string ``stop'', triggering
   fault 15. This fault sets \newline
   \verb'constant_uchar' to an out-of-domain value
   ($265$), which is wrapped back to $9$.
\end{itemize}

The results for this run are shown in Figure \ref{fig:triggers}

\begin{figure}
  \centering
  \includegraphics[scale=0.5]{graphics/verif_triggers.png}
  \caption{RUN\_triggers Results}
  \label{fig:triggers}
\end{figure}

\clearpage

\subsection{RUN\_trigger\_limited}
\textit{RUN\_trigger\_limited} is used to test the \textit{fire\_limit}
specification on the trigger-count.

The scenario is configured to apply two Overwrite-Faults, both using the same
Trigger.
The Trigger is first triggered at $t=0.2$, with a fire-limit, or
trigger-count-limit of 5.

Both Faults utilize the Trigger at $t=0.2$ and at $t=0.3$. At $t=0.4$, the
first Fault exercises the Trigger for the 5th time, and it hits its
trigger-count-limit. The second Fault sees the Trigger as non-triggered and
the Fault is not applied.

At $t=0.7$, the limiting trigger-count is removed, allowing the Fault to
reapply to both variables.

Results are confirmed in Figure
\ref{fig:trigger_limited}.

\begin{figure}[h!]
  \centering
  \includegraphics[scale=0.5]{graphics/verif_trigger_limited.png}
  \caption{RUN\_trigger\_limited Results}
  \label{fig:trigger_limited}
\end{figure}


\subsection{RUN\_errors} \label{verif_errors}
\textit{RUN\_errors} hits all the non-fatal errors in FaultManager.
Each of these errors produces an error message unless otherwise stated.
\begin{itemize}
\item Invalid trigger specification in the XML file
  \begin{itemize}
  \item Empty trigger group (no error message)
  \item Attempt to reuse a trigger that hasn't been defined earlier in the
        file (invalid reuse\_name)
  \item Missing trigger name
  \item Duplicate trigger name
  \item Missing assignment for variable
  \item Variable doesn't exist
  \item Variable is a pointer
  \item Missing comparison type (operator)
  \item Invalid comparison type (operator)
  \item Missing trigger value
  \item Invalid random value specification (\verb"RandValue")
  \item Variable is of an invalid, non-pointer type
        case)
  \item Periodic trigger with incomplete specification
  \item Periodic trigger with negative period
  \item Periodic trigger with invalid periodic length (length is greater than
        period)
  \end{itemize}
\item Invalid fault specification in XML file
  \begin{itemize}
  \item Missing fault name
  \item Duplicate fault name
  \item Missing Location
  \item Invalid Location
  \item Missing variable specification node
  \item Missing variable name
  \item Variable doesn't exist
  \item Missing fault type for Boolean
  \item Invalid fault type for Boolean
  \item Variable is of an invalid, non-pointer type
  \item Missing fault type
  \item Invalid fault type
  \item Bias faults
    \begin{itemize}
    \item Missing bias node
    \item Missing bias value
    \end{itemize}
  \item Scale faults
    \begin{itemize}
    \item Missing scale node
    \item Missing scale value
    \end{itemize}
  \item Overwrite faults
    \begin{itemize}
    \item Missing Overwrite and RandValue nodes
    \item Incomplete Overwrite and missing RandValue nodes
    \item Missing Overwrite and incomplete RandValue nodes
    \item Overwrite and RandValue nodes competing
    \end{itemize}
  \item Function faults
    \begin{itemize}
    \item Missing Function node
    \item Missing function type
    \item Invalid function type
    \item Missing independent variable node
    \item Independent variable errors
      \begin{itemize}
      \item Missing variable name
      \item Variable doesn't exist
      \item Invalid variable
      \end{itemize}
    \item Periodic functions
      \begin{itemize}
      \item Missing amplitude node
      \item Missing amplitude nominal value
      \item Missing frequency node
      \item Missing frequency nominal value
      \item Missing amplitude rate with an independent variable specified
      \item Missing amplitude independent variable with specified rate
      \end{itemize}
    \item Linear functions
      \begin{itemize}
      \item Missing parameters node
      \item Missing rate
      \end{itemize}
    \end{itemize}
  \item White-noise faults
    \begin{itemize}
    \item Missing RandValue node
    \item Invalid random number configuration
    \end{itemize}
  \item Random Walk faults
    \begin{itemize}
    \item Missing RandValue node
    \item Invalid random distribution specification
    \end{itemize}

  \item Random Number Configuration errors
    \begin{itemize}
    \item invalid random distribution specification
    \item Missing mean on Normal distribution
    \item Missing standard deviation on Normal distribution
    \item Missing minimum value on Uniform distribution
    \item Missing maximum value on Uniform distribution
    \end{itemize}

  \item Setter errors
    \begin{itemize}
    \item Invalid fault name in \textit{set\_fault\_enabled}
    \item Invalid fault name in \textit{set\_fault\_trigger\_enabled}
    \item Invalid trigger name in \textit{set\_fault\_trigger\_enabled}
    \item Invalid fault name in \textit{set\_fault\_param}
    \item Invalid trigger name in \textit{set\_trigger\_value}
    \item Invlid setting of a numeric value to a string-based trigger in
          \textit{set\_trigger\_value}
    \end{itemize}
  \end{itemize}
\end{itemize}

The faults are configured in \textit{xml\_files/errors.xml}.

This file results in four valid faults, and six valid triggers; all other
content in \textit{errors.xml} results in unusable configurations.

The four valid faults are:
\begin{itemize}
  \item \textit{Valid Fault} is used as a control, it is fully functional.
        It results in the
        application of a $+30$ bias to \verb'var_double'.
  \item \textit{Trigger Errors} has 4 valid triggers, one of which is fully
        configured and used throughout the file (\textit{valid trigger})
        and the other three are used to explore non-terminal configuration
        errors. It results in applying a bias of $+40$ to the value of
        \verb'var_float'.
  \item \textit{Redundant Spec Overwrite} evaluates a redundantly specified
        overwrite-Fault. It results in overwriting the value of
        \verb'var_float', setting the value to $2$.
  \item \textit{String Periodic Trigger} is used to evaluate the consequences
        of a downstream reassignment to the configuration of its
        \textit{string trigger} trigger. It results in the application of
        a $+2$ bias to \verb'var_double'.
\end{itemize}

There are six valid Trigger definitions:
\begin{itemize}
 \item \textit{valid trigger} is defined within the definition of the
       first Fault (\textit{Trigger Errors}). This Trigger monitors for
       the simulation time being greater than 0, which occurs on the
       second cycle. It is used for all but one of the Faults.
 \item A second Trigger named \textit{valid trigger} is defined to test
       the detection of duplicate names, which is a non-critical situation
       so it remains as a valid Trigger. This Trigger monitors for
       the simulation time being greater than 2, which is not satisfied
       during \textit{RUN\_errors}.
 \item Three Triggers are used to test faulty periodic configurations:
   \begin{itemize}
   \item \textit{Incomplete Period}
   \item \textit{Negative Period}
   \item \textit{Invalid Periodic Length}
   \end{itemize}
       Like \textit{valid trigger}, these three Triggers are also configured
       within the definition of the (\textit{Trigger Errors}) Fault and will
       trigger when simulation time is greater than 0.
       These three Triggers are configured with an invalid periodicity;
       the invalid periodicity invalidates only the periodic nature of the
       trigger, leaving each of these Triggers functionally equivalent to
       \textit{valid trigger}.
 \item \textit{string trigger} is used only in the
       \textit{String Periodic Trigger} fault; it is used to test whether a
       valid Trigger that is later given an invalid assignment will retain
       its original valid configuration.
\end{itemize}

For convenience, most (all except for \textit{string trigger}) of the
trigger configuration errors are defined within the same Fault definition
(\textit{Trigger Errors}).
There is a unique trigger defined within this fault for each of the
configuration errors.

\paragraph{Results}
The four valid faults:

\begin{itemize}
  \item \textit{Valid Fault} is configured to use the single
        \textit{valid trigger} trigger; this Fault is applied
        on the second cycle, when \textit{valid trigger} triggers.
  \item \textit{Trigger Errors} has five valid triggers (all except
        \textit{string trigger}). This Fault is
        never applied because the second \textit{valid trigger}
        Trigger does not satisfy its trigger criterion.
  \item \textit{Redundant Spec Overwrite} is configured to use the single
        \textit{valid trigger} trigger; this Fault is applied
        on the second cycle, when \textit{valid trigger} triggers.
  \item \textit{String Periodic Trigger} is configured to use the single
        \textit{string trigger} trigger; this Fault is applied
        immediately, because \textit{string trigger} is always satisfied.
\end{itemize}
The result is that:
\begin{itemize}
  \item On the first cycle, \verb'var_double' is biased by $+2$
  \item On the second cycle, \verb'var_double' is biased by $+32$
       (compound effect of \textit{String Period Trigger} and \textit{Valid
       Fault}), and \verb'var_float' is assigned the value $2$
       (\textit{Redundant Spec Override}).
\end{itemize}
These results are confirmed in Figure \ref{fig:RUN_error}.

\begin{figure}[htb]
  \centering
  \includegraphics[scale=0.5]{graphics/verif_RUN_error.png}
  \caption{Results of RUN\_error}
  \label{fig:RUN_error}
\end{figure}


The other faults defined in \textit{errors.xml} all have critical errors in
their configurations. While they all are configured to use a valid trigger,
the other errors prevent them from being made operational. None of these
other faults should cause any effect.

Examination of the simulation data (saved in verif\_data/RUN\_errors/log
\_fault\_variables.csv), the errors printed to the console, and the code
coverage report shows that this scenario functions as expected.

\subsection{Testing Terminal Errors} \label{verif_fails}
There are three scenarios that test terminal errors:
\begin{itemize}
  \item \textit{FAIL\_bad\_filename} specifies the XML filename,
        \textit{fault\_file="non-existent file"} to be a non-existent file.
        The run terminates with an error message indicating that the
        specified XML file could not be opened.
        \begin{figure}[h!]
          \centering
          \includegraphics[scale=0.5]{graphics/verif_fail_bad_filename.png}
        \end{figure}
  \item \textit{FAIL\_invalid\_xml} specifies the XML filename,
        \textit{fault\_file="xml\_files/invalid.xml"}; this file exists
        but idoes not provide the syntax necessary to be a valid XML
        definition of a \verb'FaultManager' structure hierarchy.
        \begin{figure}[h!]
          \centering
          \includegraphics[scale=0.5]{graphics/verif_fail_invalid_xml.png}
        \end{figure}
  \item \textit{FAIL\_bad\_location} attempts to execute the
        \textit{FaultManager::update(...)} method with an invalid `location'
        specification. The `location' setting indicates which set of
        Faults (e.g. \textit{UPSTREAM, DOWNSTREAM}) should be applied and
        this test attempts to use \textit{INVALID}, which is not one of
        the four recognized locations.
        \begin{figure}[h!]
          \centering
          \includegraphics[scale=0.5]{graphics/verif_fail_bad_location.png}
        \end{figure}
\end{itemize}


\section{SIM\_fault} \label{sec:verif:DI}
\par SIM\_fault instantiates four Faults and two Triggers to test specific
content related to making direct assignments to instantiated faults.
Direct assignments are not possible using the Fault Manager architecture.

The two Triggers are:
\begin{itemize}
  \item \textit{trigger} monitors the \textit{current\_time} variable (of
        type \textit{double});
        this Trigger is used in almost all test scenarios.
  \item \textit{trigger\_str} monitors \textit{control\_string} (of type
        \textit{std::string}); this Trigger is used only in the
        two \textit{RUN\_string\_trigger*} scenarios.
\end{itemize}

The four Faults are:
\begin{itemize}
  \item \textit{fault\_function} is a \textit{FaultFunction<double>} instance.
       It is used in:
  \begin{itemize}
    \item \textit{RUN\_disable\_linear\_*} to investigate the effect of
          disabling and later re-enabling the Fault mid-simulation.
    \item \textit{RUN\_linear} to provide a baseline for comparison with the
          \textit{RUN\_linear\_int}.
    \item \textit{RUN\_linear\_invalid\_comparator} and
          \textit{RUN\_linear\_periodic} as a convenient baseline to test
          additional capabilities.
    \item \textit{RUN\_new\_rate\_*} to investigate the effect of changing
          the linear function;
    \item \textit{RUN\_sinewave} to investigate the reshaping of a
          function-fault mid simulation.
    \item \textit{RUN\_string\_trigger*} to investigate the setting of
          trigger values when a Trigger is using a string variable.
  \end{itemize}
  \item \textit{fault\_function\_int} is a similar
        \textit{FaultFunction<double>}
        instance. Its independent variable is an integer, which leads to the
        appearance of a Fault using a discrete variable. This is used in
        \textit{RUN\_linear\_int}.
  \item \textit{fault\_noise} is a \textit{FaultWhiteNoise<double>} instance.
        It is used to test the random number configurations. This
        Fault is used in the four \textit{RUN\_noise\_*} runs.
  \item \textit{fault\_stale} is a \textit{FaultStale<double>} instance.
        It is used in \textit{RUN\_disable\_stale\_*} to investigate the effect
        of disabling and later re-enabling the Fault mid-simulation.
\end{itemize}

All four Faults use the same target variable.
Except in \textit{RUN\_disable\_stale\_*}, this target variable is overwritten
with a constant value at the beginning of each timestep (this is similar to the
test cases in \textit{SIM\_fault\_manager}). In \textit{RUN\_disable\_stale\_*},
the fault variable is instead overwritten with the current time at the beginning
of each timestep because overriding with a constant value makes it difficult
to discern whether the constant value is an effect of the stale Fault, or the
system override.


\subsection{Linear Function} \label{sec:verif:DI:function:linear}
In \textit{RUN\_linear\_*}, the function fault is enabled and configured to
use the linear function \newline  $var = 2 - 0.4t$.  The Trigger is
configured to apply the Fault for $t >= 1$.

\textit{RUN\_linear} produces the expected linear function after the fault
is triggered.

\textit{RUN\_linear\_int} uses \textit{time\_int}, an integer truncation of
\textit{current\_time} for its independent variable. The fault variable should
look like a set of stairs instead of a straight line.

\textit{RUN\_linear\_periodic} sets the trigger to be periodic, with the trigger
automatically oscillating between its triggered and non-triggered states. At
$t=3.5$ the periodicity is removed from the trigger and the fault applies
continually thereafter.

\textit{RUN\_linear\_invalid\_comparator} sets the triggers \textit{Operator} to
an invalid enumeration. This means the logical test for triggering will never
pass and the fault will never be applied.

Note -- the ability to adjust the periodicity of a fault is not available via
the fault manager.

The results of these runs are shown in Figure \ref{fig:linear}.

\begin{figure}[htb]
  \centering
  \includegraphics[scale=0.45]{graphics/verif_linear.png}
  \caption{Results of Linear Function Fault}
  \label{fig:linear}
\end{figure}

\subsection{Random Noise} \label{sec:verif:DI:noise}
\par In \textit{RUN\_noise\_*}, the white-noise fault is enabled and configured
to use one of the distributions outlined below. The faulted variable has a
nominal value $2$.
\begin{itemize}
 \item \textit{gaussian\_narrow}: Normal with mean 3 and standard deviation
        1, resulting in a narrow range centered on $3+2=5$ with a3-sigma
        distribution range of $[2,8]$.
 \item \textit{gaussian\_wide}: Normal with mean 3 and standard deviation 10,
        resulting in a wider range centered on $3+2=5$ with a 3-sigma
        distribution range of $[-25,35]$.
 \item \textit{uniform\_narrow}: Uniform distribution on $[-2,0]$, resulting in
       output in range $[0,2]$
 \item \textit{uniform\_wide}: Uniform distribution on $[0,10]$, resulting in
       output in range $[2,12]$
\end{itemize}

The results are shown in Figure \ref{fig:noise}.

\begin{figure}[htb]
  \centering
  \includegraphics[scale=0.45]{graphics/verif_random.png}
  \caption{Results for random number generation}
  \label{fig:noise}
\end{figure}

\clearpage

\subsection {Setting New Rate} \label{sec:verif:DI:new_rate}
\textit{RUN\_new\_rate\_*} investigates the two options for setting a new rate
mid-sim, as discussed in Section \ref{sec:formulation:fault:params}, with the
expectation illustrated in Figure \ref{fig:Rate_change}.

\textit{RUN\_new\_rate\_discontinuous} applies a direct override to the
\textit{rate} value; this leaves the \textit{nominal} value untouched and
creates a discontinuity.

\textit{RUN\_new\_rate\_continuous} applies the same override to the
\textit{rate} value, but uses the \textit{set\_param} method with the optional
3rd argument; this results in modifying the \textit{nominal} value and leaving
the fault evolution continuous.

The results of these tests are presented in Figure \ref{fig:new_rate}.


\begin{figure}[htb]
  \centering
  \includegraphics[scale=0.5]{graphics/verif_new_rate.png}
  \caption{Results for setting new rates: \newline
  Left: Discontinuous function, matching ``point of origin'' at (1,4) \newline
  Right: Continuous function; recomputed ``point of origin''}
  \label{fig:new_rate}
\end{figure}

\clearpage

\subsection {Setting Function Parameters} \label{sec:verif:DI:sinewave}
\textit{RUN\_sinewave} executes a sinusoidal function fault with varying
amplitude, frequency, and phase-offset parameters. Table \ref{tbl:sinewave}
summarizes
the changes made, more details on these changes are found in the input file at
\textit{RUN\_sinewave/input.py}

\begin{table}[htb]
\centering
\caption{String-trigger Changes} \label{tbl:sinewave}
\begin{tabular}{|l|c|c|c|} \hline
\textbf{Time} & \textbf{parameter set} & \textbf {parameter value}\\ \hline
0 & fault activation & \\ \hline
  & frequency        & 0.25\\ \hline
  & amplitude        & 1 \\ \hline
  & phase-offset     & 0.5 \\ \hline
  & amplitude-rate   & 0.2 \\ \hline
2 & phase-offset     & 0 \\ \hline
7 & amplitude-rate   & 0 \\ \hline
11 & amplitude-rate   & 0.1 \\ \hline
13 & amplitude-rate   & 0 \\ \hline
16 & amplitude        & 1 \\ \hline
   & frequency        & 0.5 \\ \hline
19 & amplitude        & -1.0 \\ \hline
21 & frequency        & 0.25 \\ \hline
   & amplitude        & 1.5 \\ \hline
24 & phase-offset-rate& 1/16 \\ \hline
   & amplitude        & 1.0 \\ \hline
28 & phase-offset-rate& 0 \\ \hline
   & amplitude        & 1.5 \\ \hline
32 & phase-offset-rate& 0.25 \\ \hline
   & amplitude        & 1.0 \\ \hline
34 & phase-offset-rate& 0 \\ \hline
   & amplitude        & 0.5 \\ \hline
35 & frequency        & 10.4 \\ \hline
   & frequency-rate   & -0.3 \\ \hline
42 & frequency-rate   & 0 \\ \hline
   & amplitude        & 0.2 \\ \hline
\end{tabular}
\end{table}

The results of these tests are presented in Figure \ref{fig:sinewave}. The
red circles in this plot identify the points at which the functional behavior
is changed. A description of these key features follows:

\begin{figure}[h!]
  \centering
  \includegraphics[scale=0.5]{graphics/verif_sinewave.png}
  \caption{Results for modifying parameters on-the-fly:
  the red dots mark the end of each segment where changes are made to
  some parameter setting}
  \label{fig:sinewave}
\end{figure}

\begin{itemize}
 \item $0<t<2$ Fault starts with phase = $0.5$, amplitude grows.
 \item $t=2$ phase-offset change results in a half-wavelength
       phase-shift. Phase jumps from $1.0$ to $0.5$. Both of these are
       zero-crossings, but the slope changes from positive to negative.
 \item $2<t<7$ Fault continues to evolve with increasing amplitude
 \item $t=7$ Amplitude-rate drops to 0 with no change to nominal amplitude.
       Total amplitude drops to $+1.0$. Phase unaffected by this change,
       remains at $1.75$ cycles, giving the Fault a value of $-1.0$ and the
       output a value of $2.0 + (-1.0) = 1.0$.
 \item $7<t<11$ Fault evolves at constant amplitude; phase advances to $2.75$
       cycles.
 \item $t=11$ Amplitude-rate set to $+0.1$; instantaneously increases the
       amplitude to $2.1$. Phase of $2.75$ cycles is a minimum, making a
       value $2.0 + (-2.1) = -0.1$.
 \item $11<t<13$ Fault evolves with increasing amplitude, growing to $2.3$ as
       the phase advances to $3.25$.
 \item $t=13$ Amplitude-rate set back to $0.0$, this time with the continuity
       flag. Amplitude remains at $2.3$. Phase = $3.25$, a maximum point so
       value = $2.0 + (+2.3) = 4.3$.
 \item $13<t<16$ Fault evolves at constant amplitude. Phase advances to
       $4.0$.
 \item $t=16$ Increase the frequency from $0.25$ to $0.5$, this should halve
       the period. Also reduce the amplitude to 1.0 to make it easier to
       identify the transition in the plot. Phase = $4.0$, a zero-crossing.
 \item $16<t<19$ Fault evolves at higher frequency and lower amplitude than
       previous. Phase advances to $5.5$.
 \item $t=19$ Amplitude is set to $-1.0$ to verify the effect of having a
       negative amplitude. Phase = $5.5$, a zero-crossing. The negation of
       the amplitude should result in the slope changing from negative (which
       is expected of a sine-wave with phase of $5.5$) to positive.
 \item $19<t<21$ The negative amplitude results in a reflection about fault=0
       and the fault continues to evolve at the same frequency and amplitude.
       Phase advances to $6.5$
 \item $t=21$ Frequency is reduced back to 0.25. Amplitude is returned to a
       positive value ($+1.5$).
       Phase = $6.5$, a zero-crossing but slope changes sign
       again, from positive to the more conventional negative associated with
       a half-phase.
 \item $21<t<24$ Fault evolves at a lower frequency again, with amplitude
       $1.5$. Phase advances to $7.25$.
 \item $t=24$ Phase = $7.25$, a maximum, so value = $2.0 + (+1.5) = 3.5$
       A phase-offset-rate is introduced, with value $1/16 = 0.0625$. The
       phase-shift assumes this phase-offset-rate has been applied since the
       start; with $dt = 24$, this creates an additional phase-offset of $1.5$,
       Phase jumps from $7.25$ to $8.75$, a maximum to a minimum.
       Amplitude is changed to 1.0. Value changes instantaneously from $3.5$
       to $2.0 + (-1.0) = 1.0$.
 \item $24<t<28$ With phase-offset-rate increasing phase and phase still
       advancing due to the integration of frequency, the apparent frequency
       increases from $0.25$ to $0.25 + 0.0625 = 0.3125$ or $5/16$. In the
       $dt=4$ of this interval, the phase should advance $5/4$ or $1.25$ from
       $8.75$ to $10.0$.
 \item $t=28$ The phase-offset-rate is set back to 0 without changing the
       nominal phase-offset, so the surplus phase resulting from the
       phase-offset-rate is removed instantaneously. The phase drops from
       $10.0$ to $8.25$, a zero-crossing to a maximum. The amplitude is
       increased again to 1.5. The value becomes $2.0 + (+1.5) = 3.5$
 \item $28<t<32$ Fault evolves at the lower frequency ($0.25$) again,
       similar to the period $21<t<24$. Phase advances to $9.25$
 \item $t=32$ The phase-offset-rate is again increased, this time to $0.25$.
       With $dt=32$, this creates a phase jump of $8.0$; the phase jumps
       from $9.25$ to $17.25$ (both represent a maximum).
       The amplitude is set back to $1.0$ to distinguish the next segment
       from the previous, causing the value to change from $2.0 + (1.5) =
       3.5$ to $2.0 + (+1.0) = 3.0$.
 \item $32<t<34$ The Fault evolves to phase $18.25$, with amplitude $1.0$
 \item $t=34$ The phase-offset-rate is returned to 0, but this time with the
       continuity flag set. The nominal phase-offset is adjusted to match
       the current phase of $18.25$. The amplitude is reduced to $0.5$
       to distinguish the next segment from the previous.
       The value to changes from $2.0 + (+1.0) = 3.0$ to $2.0 + (+0.5) = 2.5$.
 \item $34<t<35$ Fault evolves to phase 18.5.
 \item $t=35$ The final parameter to test is frequency. For the next section,
       we investigate the effect of a negative frequency. The frequency-rate
       is set to $-0.3$ and the nominal frequency is raised to $10.4$,
       leaving a net frequency of $10.4 + 35(-0.3) = -0.1$.
       The phase will now start to ``unwind'' due to the negative frequency,
       and do so at an increasing rate as the magnitude of the frequency
       increases (frequency becomes more negative). The Fault-evolution
       effectively reverses direction.
 \item $35<t<42$ As the frequency continues to get more negative, the period of
       each cycle decreases as the frequency evolves to $10.4 + 42(-0.3) =
       -2.2$.
 \item $t=42$ The frequency-rate is returned to 0 with the continuity flag set
       so the nominal frequency is adjusted to match the evolved frequency
       value of $-2.2$ from the linear function describing the frequency. The
       amplitude is reduced to $0.2$ to keep the high-frequncy oscillations
       easily observable on the plot.
 \item $t>42$ The plot confirms that the Fault continues to oscillate a little
       faster than 2Hz.
\end{itemize}


\clearpage


\subsection {Disable and Re-enable} \label{sec:verif:DI:disable}
\textit{RUN\_disable\_*} investigates the two options for disabling a fault. If
the fault is to remain disabled, the two options produce the same result, but
differ if the model is later re-enabled. This effect is discussed in section
\ref{sec:UG:DirectInstantiation:Config:Enabled}

To investigate this effect, we try two different faults -- a linear function
fault, and a stale fault. For each fault, we try two test cases. In the first,
we disable directly, setting \textit{enabled = false}, and in the second, we
disable indirectly by calling \textit{disable()}. For all four test cases, the
fault is later re-enabled directly, setting \textit{enabled = true}.

By using the \textit{disable()} method, when the fault is re-enabled, it will
reset and re-compute the initial values. Without the reset, the fault will
simply continue to use its previous initial values.

\begin{itemize}
 \item linear, direct: after re-enabling at $t=3.5$, the fault continues to
   follow the previous line. With a Fault configured to have a nominal value
   of $2.0$ and a rate of $0.4$, at $t=3.5$, the Fault-value should be $2.0 +
   0.4 (3.5 - 1.0) = 3.0$. The resulting value should be the sum of the
   nominal value of \textit{var} and the Fault-value, or $2.0 + 3.0 = 5.0$.
 \item linear, indirect: after re-enabling, the fault restarts its independent
   variable, starting back at $0.0 at t=3.5$. The output value therefore starts to
   evolve from the Fault's nominal value of $2.0$, which is added to the
   output-variable's nominal value of $2.0$ to produce an initial faulted
   value of $var=4.0$ as it did at $t=1$.  The new line is parallel to the old
   line because the \textit{rate} has not changed.
 \item stale, direct: after re-enabling, the fault returns the variable to its
   previous stale value of $1.0$.
 \item stale, indirect: after re-enabling, the fault resets the stale value to
   the current value and holds there.
\end{itemize}
Note -- recall from earlier presentation that the non-faulted value of $var$ is
constant (at $2.0$) for the linear cases, and equal to $t$ for the stale cases.

The results of these tests are presented in Figure \ref{fig:disable_enable}.

\begin{figure}
  \centering
  \includegraphics[scale=0.45]{graphics/verif_disable.png}
  \caption{Results for Disable / Re-enable}
  \label{fig:disable_enable}
\end{figure}

\clearpage

\subsection{String Trigger Behavior} \label{sec:verif:DI:string_trigger}
In \textit{RUN\_string\_trigger}, we use the \textit{trigger\_str} Trigger
to trigger the \textit{fault\_function} Fault, to fault the value of
\textit{test\_object.var} when \textit{control\_string} is equal to some
trigger-value. In this test, we want to investigate the ability to change the
trigger-value.

The ability to change trigger-values when using the Fault Manager (XML format)
implementation is limited to numeric values; the Fault Manager does not support
changing string trigger-values. This can only be accomplished when using
an independently-instantiated Trigger and Fault.

The Fault is configured to provide a linear function with a nominal value of
$2.0$ and rate of $0.4$. The sequence and expected results for this test are
presented in Table \ref{tbl:str_trigger}.
The run output data are presented in Figure \ref{fig:str_trigger}.

\begin{table}[htb]
\centering
\caption{String-trigger Changes} \label{tbl:str_trigger}
\begin{tabular}{|l|c|c|c|c|} \hline
\textbf{Time} & \textbf{control-string} & \textbf{trigger-value} &
\textbf{fault-status} & \textbf{var}\\ \hline
0 & `stop' & `go'   & off ($go \neq stop$) & $2.0$ \\ \hline
1 &        &        & off ($go \neq stop$) & $2.0$ \\ \hline
2 & `go'   &        & on  ($go = go$)      & $2.0 + (2.0 + 0.4(t-2))$ \\ \hline
3 &        & `stop' & off ($stop \neq go$) & $2.0$ \\ \hline
4 & `stop' &        & on  ($stop = stop$)  & $2.0 + (2.0 + 0.4(t-4))$ \\ \hline
\end{tabular}
\end{table}


\begin{figure}[h!]
  \centering
  \includegraphics[scale=0.5]{graphics/verif_string_trigger.png}
  \caption{Results for String Trigger Test}
  \label{fig:str_trigger}
\end{figure}

Additionally, \textit{RUN\_string\_trigger\_GT} operates the same fault with the
operator or logical comparator changed from ``is-equal-to'' to
``is-greater-than''. Evaluating relative/comparative values of strings is not
supported (evaluating whether `stop' > `go' has no unambiguous result);
this Fault never triggers and the data for this run (not shown) flatlines at
$2.0$ with no Fault applied.

\subsection{Error Cases}
This simulation provides three error cases:
\begin{itemize}
  \item \textit{ERROR\_double\_initialization} detects duplicate initialization
        on an independent variable. Because the initialization of an independent
        variable requires allocation of memory for the \verb'UntypedVariable'
        templated interface-variable (see Section
        \ref{sec:formulation:indep_var}), it should be run only once.

        Note -  only the direct-instantiation architecture requires the direct
        call to \textit{initialize()}.  The Fault Manager architecture makes
        this call behind the scenes.
  \item \textit{ERROR\_early\_initialization} detects another initialization
        problem, this time in the sequencing. To initialize a fault-function, it
        is first necessary to initialize the independent variable to allocate
        the memory for the \verb'UntypedVariable' templated interface-variable
        (see Section \ref{sec:formulation:indep_var}).

        Note -  only the direct-instantiation architecture requires the direct
        call to \textit{initialize()}.  The Fault Manager architecture makes
        this call behind the scenes.
  \item \textit{ERROR\_invalid\_function\_type} detects the invalid
        specification of the type of function to be used in a Fault-function.

        Note -  only the direct-instantiation architecture requires the direct
        setting of the function type. An invalid setting while using the Fault
        Manager architecture is detected in parsing the XML file.
\end{itemize}


\end{document}
